\documentclass[11pt,reqno]{amsart}

% ================================================================
% Paper E: E_n-Chiral Algebras, Derived Centres,
%          and the Operadic Circle
% STATUS: SECTIONS 1--5 FILLED (algebraic-geometric foundation)
% TARGET: Inventiones Mathematicae (or Annals)
% LENGTH: 80-100pp
% ================================================================

\usepackage{amsmath,amssymb,amsthm}
\usepackage{mathrsfs}
\usepackage[shortlabels]{enumitem}
\usepackage{microtype}
\usepackage{xcolor}
\usepackage[colorlinks=true,linkcolor=blue!60!black,citecolor=green!40!black]{hyperref}

\newtheorem{theorem}{Theorem}[section]
\newtheorem{proposition}[theorem]{Proposition}
\newtheorem{lemma}[theorem]{Lemma}
\newtheorem{corollary}[theorem]{Corollary}
\newtheorem{conjecture}[theorem]{Conjecture}
\theoremstyle{definition}
\newtheorem{definition}[theorem]{Definition}
\newtheorem{construction}[theorem]{Construction}
\newtheorem{example}[theorem]{Example}
\newtheorem{warning}[theorem]{Warning}
\theoremstyle{remark}
\newtheorem{remark}[theorem]{Remark}
\newtheorem{convention}[theorem]{Convention}
\newtheorem{principle}[theorem]{Principle}

% Macros
\providecommand{\cA}{\mathcal{A}}
\providecommand{\cC}{\mathcal{C}}
\providecommand{\cD}{\mathcal{D}}
\providecommand{\cO}{\mathcal{O}}
\providecommand{\cM}{\mathcal{M}}
\providecommand{\cW}{\mathcal{W}}
\providecommand{\cZ}{\mathcal{Z}}
\providecommand{\fg}{\mathfrak{g}}
\providecommand{\fh}{\mathfrak{h}}
\providecommand{\fz}{\mathfrak{z}}
\providecommand{\FM}{\overline{\mathrm{C}}}
\providecommand{\Ran}{\mathrm{Ran}}
\providecommand{\Conf}{\mathrm{Conf}}
\providecommand{\Ainf}{A_\infty}
\providecommand{\Linf}{L_\infty}
\providecommand{\Eone}{E_1}
\providecommand{\Etwo}{E_2}
\providecommand{\Ethree}{E_3}
\providecommand{\En}{E_n}
\providecommand{\Einf}{E_\infty}
\providecommand{\Pn}{\mathsf{P}_n}
\providecommand{\Pthree}{\mathsf{P}_3}
\providecommand{\SCchtop}{\mathsf{SC}^{\mathrm{ch,top}}}
\providecommand{\barB}{\bar{B}}
\providecommand{\Barord}{\bar{B}^{\mathrm{ord}}}
\providecommand{\BarSig}{\bar{B}^{\Sigma}}
\providecommand{\ChirHoch}{\mathrm{ChirHoch}}
\providecommand{\HH}{\mathrm{HH}}
\providecommand{\THH}{\mathrm{THH}}
\providecommand{\bC}{\mathbb{C}}
\providecommand{\bR}{\mathbb{R}}
\providecommand{\bQ}{\mathbb{Q}}
\providecommand{\bZ}{\mathbb{Z}}
\providecommand{\id}{\mathrm{id}}
\providecommand{\Hom}{\mathrm{Hom}}
\providecommand{\RHom}{R\!\mathrm{Hom}}
\providecommand{\Conv}{\mathrm{Conv}}
\providecommand{\End}{\mathrm{End}}
\providecommand{\MC}{\mathrm{MC}}
\providecommand{\Res}{\mathrm{Res}}
\providecommand{\gr}{\mathrm{gr}}
\providecommand{\GRT}{\mathrm{GRT}}
\providecommand{\Sym}{\mathrm{Sym}}
\providecommand{\chirAss}{\mathrm{Ass}^{\mathrm{ch}}}
\providecommand{\Spec}{\mathrm{Spec}}
\providecommand{\cF}{\mathcal{F}}
\providecommand{\cR}{\mathcal{R}}

\begin{document}

\title[$\En$-Chiral Algebras and the Operadic Circle]
{$\En$-Chiral Algebras, Derived Centres,\\
and the Operadic Circle}

\author{Raeez Lorgat}
\address{Perimeter Institute for Theoretical Physics,
 Waterloo, ON N2L 2Y5, Canada}
\email{rlorgat@perimeterinstitute.ca}

\date{\today}

\begin{abstract}
$\En$-chiral algebras are factorisation algebras on configuration
spaces of algebraic curves, inherently geometric. We construct bar
complexes at every operadic level ($n = 1, 2, 3, \infty$) on every
curve geometry (formal disk, punctured disk, annulus, nodal curves,
pair of pants), and prove that the derived chiral centre
$\cZ^{\mathrm{der}}_{\mathrm{ch}}(\cA)$ carries $\Etwo$ structure from
$\Eone$ input (Deligne) and $\Ethree$ from $\Etwo$ input (higher Deligne).
The Swiss-cheese operad $\SCchtop$ governs the pair
$(\cZ^{\mathrm{der}}_{\mathrm{ch}}(\cA), \cA)$, not the bar complex
itself. The topologisation theorem establishes
$\SCchtop + \text{conformal vector} \simeq \Ethree^{\mathrm{top}}$ for
affine Kac--Moody at non-critical level. The $\En$ operadic circle
$\Ethree \to \Etwo \to \Eone \to \Etwo \to \Ethree$ is proved to close for
simple $\fg$ via the $\Ethree$ identification theorem. Five notions of
$\Eone$-chiral algebra are distinguished. The chiral quantum group
equivalence identifies three structures ($R$-matrix, $\Ainf$,
coproduct) on the Koszul locus.
\end{abstract}

\maketitle
\tableofcontents

% ================================================================
% SECTION 1: E_n-chiral algebras on configuration spaces
% ================================================================

\section{$\En$-chiral algebras on configuration spaces}
\label{sec:en-chiral}

\subsection{The problem}

A chiral algebra is a factorisation algebra on a complex curve.
The point of departure is the collision of labelled points:
when $z_1 \to z_2$ on a curve $X$, the OPE extracts algebraic
data from the singularity along the diagonal, and the bar
construction records the iterated collisions as a coalgebra.
The bar differential uses the logarithmic $1$-form
$\eta_{ij} = d\log(z_i - z_j)$ as propagator, the Arnold
relation $\eta_{ij}\wedge\eta_{jk} + \eta_{jk}\wedge\eta_{ki}
+ \eta_{ki}\wedge\eta_{ij} = 0$ ensures $d^2 = 0$, and the
whole construction lives on Fulton--MacPherson compactifications
$\FM_k(X)$ of ordered configurations.

This is the $n = 2$ entry of a larger pattern. Replace the curve
by a real $n$-manifold, the logarithmic $1$-form by a closed
$(n{-}1)$-form, and the Arnold relation by its higher-dimensional
analogue (the Totaro relation): the result is $\En$-Koszul duality.
The chiral bar complex of a vertex algebra is the holomorphic
refinement of this topological construction, and the E$_n$ hierarchy
parametrises both the manifold dimension and the available algebraic
structure at each level.

Two axes must be kept separate throughout this paper. The
\emph{chirality axis} is the hierarchy
\begin{equation}\label{eq:chirality-hierarchy}
\Einf\text{-chiral}
\;\subset\;
P_\infty\text{-chiral}
\;\subset\;
\Eone\text{-chiral}
\end{equation}
of commutativity constraints on the OPE, all on a fixed curve.
The Heisenberg, affine Kac--Moody, and Virasoro algebras are
$\Einf$-chiral (local, symmetric OPE); the Yangian and
Etingof--Kazhdan quantum vertex algebras are $\Eone$-chiral
(nonlocal, ordered OPE). The \emph{dimensional axis} is the
$\En$ ladder
\begin{equation}\label{eq:en-ladder}
n = 1 \;(\text{interval}),
\quad
n = 2 \;(\text{surface/curve}),
\quad
n = 3 \;(\text{Chern--Simons}),
\quad
\ldots
\end{equation}
indexed by real manifold dimension. The chiral bar complex on
a complex curve belongs to $n = 2$; the $n = 1$ case is the
classical associative/$\Ainf$ bar construction on intervals.
An $\Eone$-chiral algebra on a curve is \emph{not} the same as
an $\Eone$-algebra in the topological sense: the former carries
holomorphic structure on a $2$-dimensional base, the latter lives
on a $1$-dimensional interval.

\subsection{Configuration spaces and their cohomology}

\begin{definition}[Configuration space]\label{def:e-config}
For a smooth manifold $M$ of real dimension $n$, the
\emph{ordered configuration space of $k$ points} is
\[
\Conf_k(M)
\;=\;
\{(x_1,\ldots,x_k) \in M^k : x_i \neq x_j
\text{ for } i \neq j\}.
\]
The \emph{unordered configuration space} is the quotient
$\Conf_k(M)/\Sigma_k$, where $\Sigma_k$ acts by permuting
the labels.
\end{definition}

\begin{theorem}[Arnold presentation; $n = 2$]
\label{thm:e-arnold}
The cohomology ring $H^*(\Conf_k(\bC);\bQ)$ is the
graded-commutative algebra generated by classes
$\eta_{ij} \in H^1$, $1 \le i < j \le k$, subject to the
\emph{Arnold relations}:
\begin{equation}\label{eq:e-arnold}
\eta_{ij}\wedge\eta_{jk}
+ \eta_{jk}\wedge\eta_{ki}
+ \eta_{ki}\wedge\eta_{ij}
\;=\; 0
\qquad (i < j < k).
\end{equation}
This is the Orlik--Solomon algebra $\mathrm{OS}_k$
\textup{\cite{Arnold69}}.
\end{theorem}

\begin{theorem}[Totaro presentation; general $n$
\textup{\cite{Totaro96,Coh76}}]\label{thm:e-totaro}
For $n \ge 2$, the ring $H^*(\Conf_k(\bR^n);\bQ)$ is generated
by classes $\alpha_{ij} \in H^{n-1}$, $1 \le i \neq j \le k$,
subject to:
\begin{enumerate}[label=\textup{(\roman*)}]
\item \textup{(Symmetry)}\; $\alpha_{ij} = (-1)^n\,\alpha_{ji}$.
\item \textup{(Totaro relations)}\; $\alpha_{ij}\cdot\alpha_{jk}
 + \alpha_{jk}\cdot\alpha_{ki} + \alpha_{ki}\cdot\alpha_{ij} = 0$
 for distinct $i,j,k$.
\item \textup{(Degree)}\; $\alpha_{ij}^2 = 0$ when $n$ is odd;
 when $n$ is even, $\alpha_{ij}^2$ is determined by the Euler
 class.
\end{enumerate}
At $n = 2$, identifying $\bR^2 \cong \bC$, the Totaro relations
specialize to the Arnold relations~\eqref{eq:e-arnold}. At
$n = 1$, the configuration space $\Conf_k(\bR)$ has contractible
components indexed by orderings, and $H^0$ is the group algebra
of~$\Sigma_k$; the degree-$(n{-}1) = 0$ generators encode the
ordered structure of the associative bar complex.
\end{theorem}

\begin{remark}[Arnold at $n = 2$ as Totaro specialization]
\label{rem:e-arnold-totaro}
The generators $\eta_{ij}$ of Arnold are the Totaro generators
$\alpha_{ij}$ at $n = 2$. The following table organises the
connection between the $\En$ level, the manifold, and the bar
complex:
\begin{center}
\renewcommand{\arraystretch}{1.15}
\begin{tabular}{ccll}
$n$ & Propagator degree & Manifold & Bar complex \\
\hline
$1$ & $0$ & interval/circle & associative bar (classical) \\
$2$ & $1$ & surface/curve & chiral bar (this paper) \\
$3$ & $2$ & $3$-manifold & Chern--Simons bar \\
$n$ & $n{-}1$ & $n$-manifold & $\En$ bar
\end{tabular}
\end{center}
\end{remark}

\subsection{Fulton--MacPherson compactification}

\begin{definition}[FM compactification of $\Conf_k(\bR^n)$]
\label{def:e-fm-rn}
The \emph{Fulton--MacPherson compactification}
$\overline{\Conf}_k(\bR^n)$ is the real oriented blowup
of $(\bR^n)^k$ along all partial diagonals, ordered by reverse
inclusion. It is the closure of the image of
\[
\Conf_k(\bR^n) \;\hookrightarrow\;
(\bR^n)^k \times \prod_{i<j} S^{n-1}
\times \prod_{|S| \ge 2} [0,\infty],
\]
recording positions, relative directions
$\theta_{ij} = (x_i - x_j)/|x_i - x_j|$, and relative
distances. The result is a smooth manifold with corners.
\end{definition}

\begin{proposition}[FM operad \textup{\cite{FM94}}]
\label{prop:e-fm-operad}
The codimension-$1$ boundary strata of
$\overline{\Conf}_k(\bR^n)$ are indexed by subsets
$S \subset \{1,\ldots,k\}$ with $|S| \ge 2$, and satisfy
\[
\partial_S\overline{\Conf}_k(\bR^n)
\;\cong\;
\overline{\Conf}_{k - |S| + 1}(\bR^n)
\times
\overline{\Conf}_{|S|}(\bR^n).
\]
This factorisation equips
$\{\overline{\Conf}_k(\bR^n)\}_{k \ge 0}$ with the structure
of the FM operad $\mathsf{FM}_n$, a model for $\En$
\textup{\cite{GJ,SalvatoreLongoni05}}.
\end{proposition}

\subsection{The $\En$ propagator}

\begin{definition}[$\En$ propagator]\label{def:e-propagator}
An \emph{$\En$ propagator} on an oriented $n$-manifold $M$ is
a smooth closed $(n{-}1)$-form
$G \in \Omega^{n-1}(M \times M \setminus \Delta)$ satisfying:
\begin{enumerate}[label=\textup{(\roman*)}]
\item \textup{(Poincar\'e dual of diagonal)}\;
 $\int_{M \times M} G \wedge \pi_2^*\eta = \int_M \eta$
 for every closed $\eta$.
\item \textup{(Singularity)}\;
 Near $\Delta$, in geodesic coordinates $r = x_1 - x_2$,
 \[
 G \;\sim\;
 \frac{1}{\mathrm{vol}(S^{n-1})}\,
 \frac{\sum_{a=1}^n (-1)^{a-1} r_a\,
 dr_1 \wedge \cdots \wedge \widehat{dr_a}
 \wedge \cdots \wedge dr_n}{|r|^n}
 \;+\;\text{smooth}.
 \]
\item \textup{(Symmetry)}\;
 $\sigma^*G = (-1)^n\,G$ under the swap
 $(x_1,x_2) \mapsto (x_2,x_1)$.
\end{enumerate}
\end{definition}

\begin{example}[$n = 2$: the chiral propagator]
\label{ex:e-n2-chiral}
For $M = \bC$, the propagator is $G = \eta_{12} = d\log(z_1 - z_2)$,
a meromorphic $1$-form. Residues along $z_1 = z_2$ extract OPE data.
The symmetry $\sigma^*G = d\log(z_2 - z_1) = G$ holds since
$d\log(-1) = 0$.
\end{example}

\begin{example}[$n = 3$: the Chern--Simons propagator]
\label{ex:e-n3-cs}
For $M = \bR^3$, the propagator is the closed $2$-form
\[
G \;=\; \frac{1}{4\pi}\,
\frac{\epsilon_{abc}\, r^a\, dr^b \wedge dr^c}{|r|^3},
\qquad r = x_1 - x_2.
\]
Integration over $S^2$ gives the linking number:
$\int_{S^2} G = 1$.
\end{example}

\begin{example}[$n = 1$: the associative case]
\label{ex:e-n1-assoc}
For $M = \bR$, the ``propagator'' is a $0$-form: the sign function
$G(x_1,x_2) = \mathrm{sgn}(x_1 - x_2)$, which records the ordering
of points on the real line. The bar complex is the classical tensor
coalgebra $T^c(s^{-1}\bar{A})$ with deconcatenation coproduct;
there is no de~Rham differential to speak of, and the Arnold/Totaro
relations are trivial (degree-$0$ generators commute vacuously).
\end{example}

\begin{proposition}[Residue as linking sphere integral]
\label{prop:e-linking-residue}
For an $\En$ propagator $G$ and a smooth function $f$ near the
diagonal, the $\En$ residue is
\[
\Res_\Delta(G \wedge f)
\;:=\;
\lim_{\varepsilon \to 0}
\int_{|x_1 - x_2| = \varepsilon} G \cdot f(x_1,x_2)
\;=\; f(x,x).
\]
This generalises the chiral residue $\Res_{z_1 = z_2}$ at $n = 2$.
\end{proposition}

\begin{proof}
The singularity condition forces $G|_{|r| = \varepsilon}$ to
be the normalised volume form on $S^{n-1}$ up to terms vanishing
as $\varepsilon \to 0$. The integral therefore extracts $f(x,x)$.
\end{proof}

\subsection{The $\En$ bar complex}

\begin{definition}[$\En$ bar complex]\label{def:e-en-bar}
Let $A$ be an $\En$-algebra. The \emph{$\En$ bar complex} is the
chain complex
\begin{equation}\label{eq:e-en-bar}
\barB_{\En}(A)
\;=\;
\bigoplus_{k \ge 1}
\Omega^*\!\bigl(\overline{\Conf}_k(\bR^n)\bigr)
\otimes_{\Sigma_k} (s^{-1}\bar{A})^{\otimes k},
\end{equation}
with differential
\begin{equation}\label{eq:e-en-bar-diff}
d_{\En}
\;=\; d_{\mathrm{dR}}
\;+\;
\sum_{\substack{S \subset \{1,\ldots,k\} \\ |S| \ge 2}}
\Res_{\partial_S}\!(G_S \wedge \mu_S),
\end{equation}
where $G_S$ is the product of propagators for pairs within $S$,
$\Res_{\partial_S}$ is the linking-sphere residue along the
boundary stratum $\partial_S$ of $\overline{\Conf}_k(\bR^n)$,
and $\mu_S \colon \bar{A}^{\otimes|S|} \to \bar{A}$ is the
$\En$-multiplication.
\end{definition}

\begin{theorem}[$d^2 = 0$ from Totaro relations]
\label{thm:e-d2-totaro}
The differential $d_{\En}$ satisfies $d_{\En}^2 = 0$.
\end{theorem}

\begin{proof}
Decompose $d_{\En}^2 = d_{\mathrm{dR}}^2
+ [d_{\mathrm{dR}}, d_{\mathrm{res}}] + d_{\mathrm{res}}^2$.
The first term vanishes. The cross-term vanishes because $G$ is
closed. The square $d_{\mathrm{res}}^2$ involves iterated
residues: for three points $i,j,k$, the three iterated residues
cancel by the Totaro relation
$\alpha_{ij}\cdot\alpha_{jk} + \alpha_{jk}\cdot\alpha_{ki}
+ \alpha_{ki}\cdot\alpha_{ij} = 0$. Nested collisions cancel by
the $\En$-algebra associativity coherences encoded in the FM
operad structure (Proposition~\ref{prop:e-fm-operad}).
\end{proof}

\begin{remark}[$n$-dependence of the bar complex]
\label{rem:e-n-dependence}
At $n = 1$: the propagator has degree $0$, the de~Rham
differential vanishes, and the bar complex reduces to
$T^c(s^{-1}\bar{A})$ with the classical bar differential from
the algebra multiplication. At $n = 2$: the propagator is a
$1$-form, and the bar complex is the chiral bar complex
$\barB^{\mathrm{ch}}(\cA)$ of this paper. At $n = 3$: the
propagator is a $2$-form, and the bar complex computes
perturbative Chern--Simons invariants.
\end{remark}

\subsection{The four levels}

The operadic levels relevant to chiral algebras on curves are:

\begin{definition}[Four operadic levels on a curve]
\label{def:e-four-levels}
Let $X$ be a smooth complex curve. We distinguish four levels
of algebraic structure:
\begin{enumerate}[label=\textup{(\roman*)}]
\item \emph{$\Einf$-chiral} (local, symmetric OPE):
 the factorisation algebra is defined on unordered
 configuration spaces, with $\Sigma_k$-equivariant
 multiplication.
 All standard vertex algebras (Heisenberg, affine Kac--Moody,
 Virasoro, $\cW_N$, free-field algebras) are $\Einf$-chiral.
\item \emph{$\Etwo$-topological}: the bar complex
 $\barB_{\Etwo}(A)$ uses the topological propagator
 $(1/2\pi)\,d\arg(z_1 - z_2)$ and depends only on the
 oriented topology of~$X$.
\item \emph{$\Eone$-chiral} (ordered, nonlocal OPE):
 the factorisation algebra lives on ordered configuration
 spaces $\Conf_k^{\mathrm{ord}}(X)$, with no $\Sigma_k$
 quotient. Yangians and Etingof--Kazhdan quantum vertex
 algebras are $\Eone$-chiral.
\item \emph{$\Ethree$-topological}: the algebra of bulk
 observables. Accessed via the derived centre
 (topologisation theorem, Section~\ref{sec:topologisation}).
\end{enumerate}
\end{definition}

\begin{remark}[The chirality axis is orthogonal to the
dimensional axis]\label{rem:e-orthogonality}
An $\Einf$-chiral algebra and an $\Eone$-chiral algebra are both
``$n = 2$ objects'' in the dimensional ladder: they are both
factorisation algebras on a complex curve. The distinction is
the commutativity of the OPE, not the manifold dimension.
The $\Etwo$-topological bar complex uses $n = 2$ topological
configurations; the $\Eone$ ordered bar complex uses the
$\Eone$-chiral ordering. These are different structures on the
same underlying space.
\end{remark}

\subsection{The Heisenberg at each level}

The Heisenberg algebra $H_k$ at level $k$, generated by a
single field $J(z)$ with OPE $J(z)J(w) \sim k/(z-w)^2$,
provides the universal test case. As an abelian chiral algebra,
its structure simplifies at every operadic level.

\begin{example}[Heisenberg: $\Einf$-chiral]
\label{ex:e-heis-einf}
The Heisenberg $H_k$ is $\Einf$-chiral. The OPE has a double
pole $k/(z-w)^2$ and no simple pole (the algebra is abelian),
so all OPE products are $\Sigma_k$-symmetric. The
$\Einf$-chiral bar complex is the coshuffle coalgebra
$S^c(s^{-1}\overline{H}_k)$ with the bar differential encoding
the double-pole OPE. The classical $r$-matrix is
% AP126 check: r(z)|_{k=0} = 0 (abelian limit). Y.
\begin{equation}\label{eq:e-heis-r}
r^{\mathrm{Heis}}(z) \;=\; k/z,
\end{equation}
with $\kappa(H_k) = \mathrm{av}(r(z)) = k$.
At $k = 0$, $r(z) = 0$ and $\kappa = 0$: the algebra becomes
trivially abelian with no scattering.
\end{example}

\begin{example}[Heisenberg: $\Eone$-chiral]
\label{ex:e-heis-e1}
The ordered bar complex
$\Barord(H_k) = T^c(s^{-1}\overline{H}_k)$ has deconcatenation
coproduct. Since $H_k$ is $\Einf$, the formality bridge holds:
$\Barord(H_k) \xrightarrow{\sim} \BarSig(H_k)$ is a
quasi-isomorphism. The $R$-matrix is trivial (the Heisenberg
is commutative), and the averaging map
$\mathrm{av}\colon\fg^{\Eone}_{H_k} \to \fg^{\mathrm{mod}}_{H_k}$
loses nothing. The Heisenberg is class~G (shadow tower
terminates at degree $2$: $\Delta = 8\kappa\,S_4 = 0$
because $S_4 = 0$).
\end{example}

\begin{example}[Heisenberg: $\Etwo$-topological]
\label{ex:e-heis-e2}
The $\Etwo$ bar complex
$\barB_{\Etwo}(H_k^{\mathrm{top}})$ uses the topological
propagator $(1/2\pi)\,d\arg(z_1 - z_2)$. The bar cohomology
is concentrated in degree $1$: the Heisenberg is Koszul. The
Koszul dual is $H_k^! = \Sym^{\mathrm{ch}}(V^*)$, the
symmetric chiral algebra on the dual generator. The Koszul
complementarity is
$\kappa(H_k) + \kappa(H_k^!) = k + (-k) = 0$.
\end{example}

\begin{example}[Heisenberg: $\Ethree$-topological]
\label{ex:e-heis-e3}
The derived chiral centre
$\cZ^{\mathrm{der}}_{\mathrm{ch}}(H_k)$ carries
$\Etwo$ from the chiral direction and $\Eone$ from the
topological direction. Since $H_k$ has a conformal vector
at any $k \neq 0$ (the Sugawara element
$T = (1/2k)\,{:}JJ{:}$ at level $k$), topologisation applies:
the BRST cohomology carries $\Ethree^{\mathrm{top}}$.
For the Heisenberg, this $\Ethree$ structure is trivial
(all operations are determined by the commutative product
and the linear Lie bracket), reflecting class G.
\end{example}

\begin{remark}[The Heisenberg as Casimir--Goldstone opening]
\label{rem:e-heis-cg}
The Heisenberg is the abelian Kac--Moody algebra:
the same object as the level-$k$ $\mathrm{U}(1)$ current
algebra, as the boundary algebra of abelian Chern--Simons
theory, and as the universal target of the Casimir element
in the chiral bar complex. Every formula in this paper is
computed first for $H_k$ and then extended to general
families by variation of the $r$-matrix.
\end{remark}

% ================================================================
% SECTION 2: Bar chain models on curve geometries
% ================================================================

\section{Bar chain models on curve geometries}
\label{sec:bar-chain-models}

The $\En$ bar complex of Section~\ref{sec:en-chiral} is defined
on configuration spaces of $\bR^n$: a flat, contractible base
with no periods, no curvature, no modular anomaly. A chiral
algebra lives on curves whose global topology shapes the bar
differential. On the formal disk $D$, the bar complex is flat
($d^2 = 0$). On a genus-$g$ curve $\Sigma_g$, the period matrix
introduces curvature $d_{\mathrm{fib}}^2 = \kappa \cdot \omega_g$.
The passage from $D$ to $\Sigma_g$ is the passage from algebra
to geometry.

\subsection{The formal disk $D$}

\begin{definition}[Bar complex on the formal disk]
\label{def:e-bar-disk}
Let $D = \Spec\bC[[z]]$ be the formal disk and $\cA$ a chiral
algebra on~$D$. The \emph{bar complex on $D$} is
\begin{equation}\label{eq:e-bar-disk}
\barB_D(\cA) \;=\; T^c(s^{-1}\bar{\cA}),
\qquad
d_{\barB} = \sum_{i < j}
\Res_{z_i = z_j}\!
\bigl(\eta_{ij} \wedge \mu_{ij}\bigr),
\end{equation}
where $\eta_{ij} = d\log(z_i - z_j)$ is the chiral propagator,
$\mu_{ij}$ is the OPE product at positions $z_i, z_j$, and
$\bar{\cA} = \ker(\varepsilon)$ is the augmentation ideal.
The coproduct is deconcatenation:
\[
\Delta[a_1|\cdots|a_n]
\;=\;
\sum_{i=0}^{n} [a_1|\cdots|a_i]
\otimes [a_{i+1}|\cdots|a_n].
\]
This makes $\barB_D(\cA)$ an $\Eone$-chiral coassociative
coalgebra (a dg coalgebra over $(\chirAss)^!$).
The bar complex on the formal disk sees genus $0$ only. There
is no period matrix, no curvature, no modular anomaly. The bar
differential uses only the collision data of points on a
contractible space.
\end{definition}

\begin{proposition}[$d^2 = 0$ on $D$]
\label{prop:e-bar-disk-d2}
The differential~\eqref{eq:e-bar-disk} satisfies $d_{\barB}^2 = 0$.
\end{proposition}

\begin{proof}
By Theorem~\ref{thm:e-d2-totaro} at $n = 2$: the Arnold
relation on $\Conf_k(D)$ ensures cancellation of iterated
residues. On the formal disk, the global topology is trivial
and there is no curvature correction.
\end{proof}

\begin{example}[Heisenberg on $D$]\label{ex:e-heis-disk}
The bar complex $\barB_D(H_k)$ has generators
$s^{-1}J_{-n}\,\mathbf{1}$ for $n \ge 1$, where
$J_{-n} = (1/(n{-}1)!)\,\partial^{n-1}J$. The bar
differential encodes the OPE $J(z)J(w) \sim k/(z-w)^2$:
the only nonzero component of $d_{\barB}$ at degree $2$ is
\[
d_{\barB}[J_{-1}|J_{-1}]
\;=\; k\,[J_{-2}],
\]
the residue of the double-pole OPE contracted against
$d\log(z_1 - z_2)$, which absorbs one pole to produce a
simple-pole residue. Bar cohomology is concentrated in
degree $1$: the Heisenberg is Koszul.
\end{example}

\subsection{The punctured disk $D^*$}

\begin{definition}[Bar complex on $D^*$]
\label{def:e-bar-punc}
Let $D^* = \Spec\bC((z))$ be the punctured formal disk. The
bar complex on $D^*$ is
\begin{equation}\label{eq:e-bar-punc}
\barB_{D^*}(\cA)
\;=\;
T^c(s^{-1}\bar{\cA}),
\qquad
d_{\barB}
= d_D
+ \sum_i \oint_{|z_i| = \varepsilon} \eta_i \wedge \partial_i,
\end{equation}
where $d_D$ is the disk bar differential~\eqref{eq:e-bar-disk}
and the second term records the monodromy around the puncture.
The punctured disk is the home of the state-field
correspondence: a bar element of degree $k$ is a $k$-fold
collision pattern that wraps around the origin.
\end{definition}

\begin{proposition}[The $\Etwo$-$\Eone$ gap on $D^*$]
\label{prop:e-e2-e1-gap}
On the punctured disk, the ordered bar complex
$\Barord_{D^*}(\cA)$ carries strictly more information than the
symmetric bar $\BarSig_{D^*}(\cA)$. The averaging map
$\mathrm{av}\colon \Barord \to \BarSig$ has kernel controlled
by braid-group representations of
$\pi_1(\Conf_k(D^*)) = B_k^{\mathrm{aff}}$
\textup{(}the affine braid group\textup{)}.
\end{proposition}

\begin{proof}
The fundamental group $\pi_1(\Conf_k(D^*))$ is the affine
braid group on $k$ strands, which surjects onto $\Sigma_k$ with
kernel the pure affine braid group. The symmetric bar complex
takes $\Sigma_k$-coinvariants, killing the braid-group
representations. For $\Einf$-chiral algebras (Heisenberg,
Kac--Moody), the $R$-matrix is trivial and the loss is
inessential. For $\Eone$-chiral algebras (Yangian), the
$R$-matrix is nontrivial and the kernel of $\mathrm{av}$
contains the KZ associator and its higher-degree descendants.
\end{proof}

\begin{example}[Heisenberg on $D^*$]\label{ex:e-heis-punc}
The monodromy of $H_k$ around the puncture is trivial (the
$R$-matrix $r(z) = k/z$ has no branch cut), so
$\barB_{D^*}(H_k) \simeq \barB_D(H_k)$. The $\Etwo$-$\Eone$
gap is zero: the ordered and symmetric bars agree.
\end{example}

\subsection{The annulus}

\begin{definition}[Bar complex on the annulus]
\label{def:e-bar-annulus}
Let $A_{r,R} = \{z : r < |z| < R\}$ be an annulus. The bar
complex on $A_{r,R}$ computes the topological Hochschild homology of $\cA$
as a chiral algebra:
\begin{equation}\label{eq:e-bar-annulus}
\barB_{A_{r,R}}(\cA)
\;\simeq\;
\barB(\cA) \otimes^{\mathbf{L}}_{\cA^e} \barB(\cA),
\end{equation}
where $\cA^e = \cA \otimes \cA^{\mathrm{op}}$ is the
enveloping algebra and the derived tensor product computes
the self-Tor of the bar resolution. The annulus has two boundary
circles (inner and outer), giving it the structure of a bimodule
over~$\cA$.
\end{definition}

\begin{proposition}[Annulus computes chiral Hochschild homology]
\label{prop:e-annulus-hh}
The factorisation homology of $\cA$ over the annulus
$A_{r,R}$ is
\[
\int_{A_{r,R}} \cA
\;\simeq\;
\ChirHoch_*(\cA).
\]
The inner boundary gives the left $\cA$-module structure, the
outer boundary the right $\cA$-module structure, and the
factorisation homology of the annulus is the derived trace.
\end{proposition}

\begin{proof}
The annulus is the trace geometry for factorisation homology:
$\int_{S^1 \times I} \cA \simeq \HH_*(\cA)$ by the
K\"unneth property of factorisation homology applied to the
product $S^1 \times I$, where $S^1$ contributes the trace
and $I$ the bimodule structure. The chiral refinement replaces
the topological $S^1$ with the holomorphic annulus and the
topological HH with $\ChirHoch$.
\end{proof}

\begin{example}[Heisenberg on the annulus]\label{ex:e-heis-annulus}
The chiral Hochschild homology of $H_k$ is
$\ChirHoch_*(H_k) \cong \bC[J_{-1}, J_{-2}, \ldots]$, the
polynomial algebra on the negative modes. The annulus partition
function is the character
$\mathrm{tr}_{H_k}(q^{L_0}) = 1/\eta(q)$, where
$\eta(q) = q^{1/24}\prod_{n \ge 1}(1 - q^n)$
is the Dedekind eta function.
\end{example}

\subsection{Nodal curves and curvature}

\begin{definition}[Bar complex at genus $g$]\label{def:e-bar-genus}
Let $\Sigma_g$ be a smooth genus-$g$ curve and $\cA$ a chiral
algebra on $\Sigma_g$. The genus-$g$ bar complex is
\begin{equation}\label{eq:e-bar-genus}
\barB^{(g)}(\cA)
\;=\;
\bigoplus_{k \ge 1}
\Omega^*\!\bigl(\FM_k(\Sigma_g)\bigr)
\otimes (s^{-1}\bar{\cA})^{\otimes k},
\end{equation}
with differential $d_g = d_{\mathrm{dR}} + d_{\mathrm{res}}$,
where $d_{\mathrm{res}}$ uses the genus-$g$ propagator
$d\log E(z,w)$, with $E(z,w)$ the prime form on $\Sigma_g$.
The propagator $d\log E(z,w)$ has weight $1$ for all $g$
\textup{(}the prime form is a section of
$K^{-1/2} \boxtimes K^{-1/2}$\textup{)}.
\end{definition}

\begin{proposition}[Curvature at genus $g \ge 1$]
\label{prop:e-curvature}
At genus $g \ge 1$, the bar differential is curved:
\begin{equation}\label{eq:e-curvature}
d_{\mathrm{fib}}^2
\;=\; \kappa(\cA) \cdot \omega_g,
\end{equation}
where $\omega_g = c_1(\lambda) \in H^2(\overline{\cM}_g)$ is
the first Chern class of the Hodge line bundle on the moduli
space $\overline{\cM}_g$, and $\kappa(\cA)$ is the scalar
curvature of $\cA$. The curvature is on the fiberwise bar
differential over $\overline{\cM}_g$, not on the total bar
differential \textup{(}which satisfies $d^2 = 0$ on the
convolution algebra over $\overline{\cM}_{g,n}$\textup{)}.
\end{proposition}

\begin{proof}
The genus-$g$ propagator $d\log E(z,w)$ acquires period
corrections from $H^0(\Sigma_g, \Omega^1)$. These corrections
introduce the curvature $\kappa \cdot \omega_g$ when the bar
differential is restricted to fibers of the universal curve
over $\overline{\cM}_g$. The computation reduces to the
anomaly of the Hodge line bundle and is performed in the
convolution algebra formalism.
\end{proof}

\begin{example}[Heisenberg at genus $g$]\label{ex:e-heis-genus}
At genus $g \ge 1$, the curvature is
$d_{\mathrm{fib}}^2 = k \cdot \omega_g$, using
$\kappa(H_k) = k$. At $k = 0$, the curvature vanishes: the
bar complex is uncurved at every genus. The genus-$g$ anomaly
$\mathrm{obs}_g = \kappa \cdot \lambda_g = k \cdot \lambda_g$
(uniform-weight) gives $F_1 = k/24$ from the Hodge class on
$\overline{\cM}_{1,1}$.
\end{example}

\subsection{Pair of pants and fusion}

\begin{definition}[Pair of pants]\label{def:e-pants}
The pair of pants $P$ is a genus-$0$ surface with three boundary
circles. The factorisation homology over $P$ computes the fusion
product:
\[
\int_P \cA \;=\; \cA \otimes_{\cA}^{\mathbf{L}} \cA,
\]
where the derived tensor product is taken over the chiral algebra
itself, viewed as a bimodule via left and right OPE.
\end{definition}

\begin{proposition}[Fusion from bar complex]\label{prop:e-fusion}
The fusion product of two $\cA$-modules $M_1, M_2$ is computed
by the two-pointed bar complex:
\[
M_1 \otimes_\cA^{\mathbf{L}} M_2
\;\simeq\;
\barB(\cA;\, M_1, M_2)
\;=\;
\bigoplus_{k \ge 0}
M_1 \otimes (s^{-1}\bar{\cA})^{\otimes k} \otimes M_2,
\]
with differential contracting adjacent factors via OPE. The
associated derived category is the category of conformal blocks.
\end{proposition}

\begin{example}[Heisenberg fusion]\label{ex:e-heis-fusion}
The Fock modules $\cF_\lambda$ of $H_k$ at momentum $\lambda$
fuse by
$\cF_{\lambda_1} \otimes_{H_k}^{\mathbf{L}} \cF_{\lambda_2}
\simeq \cF_{\lambda_1 + \lambda_2}$,
reflecting the additive structure of the abelian current
algebra. The fusion rules are trivial: all Fock modules fuse
freely. This is the class~G triviality.
\end{example}

% ================================================================
% SECTION 3: The derived chiral centre
% ================================================================

\section{The derived chiral centre}
\label{sec:derived-centre}

\subsection{The problem}

Given a chiral algebra $\cA$, what is the algebra of bulk
observables of the associated $3$-dimensional
holomorphic--topological theory? The answer is the derived
chiral centre
$\cZ^{\mathrm{der}}_{\mathrm{ch}}(\cA) =
C^\bullet_{\mathrm{ch}}(\cA, \cA)$, the chiral Hochschild
cochain complex. This is not the commutant of $\cA$ in itself
(the naive centre, which may be too small), but the derived
endomorphism object computed via the bar resolution.

The key structural result is that this derived centre carries
$\Etwo$ structure from the chiral direction (the chiral Deligne
conjecture, Theorem~\ref{thm:e-center-hochschild} below) and
$\Eone$ structure from the topological direction, giving a
total $\SCchtop$ structure on the \emph{pair}
$(\cZ^{\mathrm{der}}_{\mathrm{ch}}(\cA), \cA)$.

\subsection{Chiral Hochschild cochains}

\begin{definition}[Chiral Hochschild cochain complex]
\label{def:e-chirHoch}
Let $\cA$ be a chiral algebra on a curve $X$ and let
$\End^{\mathrm{ch}}_\cA$ denote the chiral endomorphism
operad (the collection of multilinear maps
$\cA^{\otimes m} \to \cA$ parametrised by points on $X$ and
their collision patterns). The \emph{chiral Hochschild cochain
complex} is
\begin{equation}\label{eq:e-chirHoch}
C^\bullet_{\mathrm{ch}}(\cA, \cA)
\;=\;
\prod_{m \ge 0}
\Hom_{\Sigma_m}\!\bigl(
\FM_m(\bC),\,
\Hom(\cA^{\otimes m}, \cA)
\bigr),
\end{equation}
with differential combining the bar differential of the
configuration space with the internal differential of $\cA$.
A degree-$k$ cochain is a map that takes $m$ inputs from
$\cA$, arranged on the configuration space $\FM_m(\bC)$, and
produces a single output in $\cA$ with the correct
equivariance.
\end{definition}

\begin{definition}[Chiral braces]
\label{def:e-chiral-braces}
The cochain complex $C^\bullet_{\mathrm{ch}}(\cA, \cA)$
carries \emph{chiral brace operations}
\begin{equation}\label{eq:e-braces}
\{f; g_1, \ldots, g_k\}
\;\in\;
C^{p+q_1+\cdots+q_k - k}_{\mathrm{ch}}(\cA, \cA)
\end{equation}
for $f \in C^p$ and $g_i \in C^{q_i}$, defined by inserting
the $g_i$ into the inputs of $f$ at specified positions in
$\FM_m(\bC)$. These are the chiral analogues of the classical
Gerstenhaber braces. The brace operations satisfy:
\begin{enumerate}[label=\textup{(\roman*)}]
\item the pre-Lie identity on the single insertion
 $\{f; g\}$ (which defines the chiral Gerstenhaber bracket
 $[f, g] = \{f; g\} - (-1)^{(|f|-1)(|g|-1)}\{g; f\}$);
\item higher coherence relations making
 $C^\bullet_{\mathrm{ch}}(\cA, \cA)$ a homotopy
 $\Etwo$-algebra.
\end{enumerate}
\end{definition}

\subsection{The chiral centre theorem}

\begin{theorem}[Chiral centre theorem: the operadic centre of
$\SCchtop$ is $C^\bullet_{\mathrm{ch}}$]
\label{thm:e-center-hochschild}
Let $\cA$ be an $\Einf$-chiral algebra on a curve $X$ and
$A = \cA|_{\{z_0\}}$ the fiber at a basepoint. The
\emph{operadic centre} of $A$ in the Swiss-cheese operad
$\SCchtop$ is quasi-isomorphic to the chiral Hochschild cochain
complex as $\Etwo$-algebras:
\begin{equation}\label{eq:e-center-qi}
Z_{\SCchtop}(A)
\;\simeq\;
C^\bullet_{\mathrm{ch}}(\cA, \cA).
\end{equation}
The $\Etwo$-structure comes from the closed sector
$\{\FM_k(\bC)\}_{k \ge 0}$ of $\SCchtop$: the FM
compactification of the plane provides the operadic parameter
space for the braces and the cup product.
\end{theorem}

\begin{proof}
By the operadic centre construction
(Definition~\ref{def:e-operadic-centre} below), an element
of $Z_{\SCchtop}(A)$ at degree $(k,m)$ is a map
\[
\phi_{k,m}\colon
\FM_k(\bC) \times \Eone(m) \;\longrightarrow\;
\Hom(A^{\otimes m}, A),
\]
equivariant for $\Sigma_k$ on $\FM_k(\bC)$ and $\Sigma_m$
on $\Eone(m)$. The $\Eone(m)$ factor provides the open inputs;
the $\FM_k(\bC)$ factor provides the closed/holomorphic
inputs. Taking the limit over all $m$, the components assemble
into the chiral Hochschild cochain complex~\eqref{eq:e-chirHoch}.
The closed-sector composition maps in $\SCchtop$ (FM insertion
in~$\bC$) produce the brace operations, and the compatibility
conditions of the operadic centre are precisely the higher
coherence relations that make the braces into a homotopy
$\Etwo$-structure.
\end{proof}

\begin{definition}[Operadic centre]
\label{def:e-operadic-centre}
Let $\cO$ be a two-coloured operad with colours
$\{\mathsf{c}, \mathsf{o}\}$ and the vanishing constraint
$\cO(\ldots, \mathsf{o}, \ldots;\, \mathsf{c}) = \varnothing$
(no open-to-closed maps). For an
$\cO^{\mathsf{o}}$-algebra $A$, the \emph{operadic centre}
$Z_\cO(A)$ is the terminal object in the category of
$\cO$-algebra pairs $(B, A)$ with fixed open colour~$A$:
\[
Z_\cO(A) \;=\;
\varprojlim_{\substack{(B,A) \in \\
\cO\text{-}\mathsf{Alg}_{/A}}} B.
\]
Concretely, it is the equaliser in
\[
Z_\cO(A) \;\hookrightarrow\;
\prod_{m \ge 0}
\Hom_{\Sigma_m}\!\bigl(
\cO(\mathsf{c}, \mathsf{o}^m;\, \mathsf{o}),\,
\Hom(A^{\otimes m}, A)\bigr)
\;\rightrightarrows\; (\cdots),
\]
where the two maps enforce closed-sector and open-sector
equivariance.
\end{definition}

\subsection{Theorem H: concentration in degrees $\{0,1,2\}$}

\begin{theorem}[Theorem H: cohomological concentration]
\label{thm:e-theorem-H}
Let $\cA$ be a Koszul chiral algebra. The chiral Hochschild
cohomology $\ChirHoch^*(\cA) = H^*(C^\bullet_{\mathrm{ch}}(\cA,\cA))$
is concentrated in cohomological degrees $\{0,1,2\}$:
\begin{equation}\label{eq:e-thmH}
\ChirHoch^i(\cA) = 0
\qquad \text{for } i \ge 3.
\end{equation}
The Hilbert series is polynomial. The three nonvanishing degrees
have the following interpretations:
\begin{enumerate}[label=\textup{(\roman*)}]
\item $\ChirHoch^0(\cA) = Z(\cA)$: the naive centre.
\item $\ChirHoch^1(\cA)$: first-order deformations of $\cA$
 as a chiral algebra.
\item $\ChirHoch^2(\cA)$: obstructions to extending
 deformations.
\end{enumerate}
\end{theorem}

\begin{remark}[Amplitude, not virtual dimension]
\label{rem:e-amplitude}
Concentration in degrees $\{0,1,2\}$ is a statement about
\emph{cohomological amplitude}: the cochain complex is
quasi-isomorphic to one supported in degrees $0, 1, 2$. This
is not the same as ``virtual dimension $2$'' (an arithmetic
invariant) or ``$\ChirHoch^*(\cA) = \bC[\Theta]$''
(which is false for all non-trivial algebras). It also differs
from the Gelfand--Fuchs cohomology $H^*_{\mathrm{GF}}(\fg)$,
which is infinite-dimensional for the Virasoro algebra.
\end{remark}

\begin{example}[Heisenberg: $\ChirHoch^*(H_k)$]
\label{ex:e-heis-chirhoch}
For the Heisenberg algebra $H_k$:
\begin{align}
\ChirHoch^0(H_k) &= \bC
\qquad\text{(the centre is $1$-dimensional)},
\label{eq:e-heis-hoch0}\\
\ChirHoch^1(H_k) &= \bC
\qquad\text{(one-parameter deformation: the level $k$)},
\label{eq:e-heis-hoch1}\\
\ChirHoch^2(H_k) &= \bC
\qquad\text{(one obstruction class)}.
\label{eq:e-heis-hoch2}
\end{align}
Total dimension: $\dim\ChirHoch^*(H_k) = 3$.
For comparison: the naive centre $Z(H_k) = \bC$ has dimension
$1$; the derived centre has dimension $3$, detecting the
$\mathrm{Ext}^1$ and $\mathrm{Ext}^2$ contributions invisible
to the commutant.
\end{example}

\begin{example}[Affine Kac--Moody: $\ChirHoch^*(V_k(\fg))$]
\label{ex:e-km-chirhoch}
For the affine Kac--Moody algebra $V_k(\fg)$ with
$\fg$ simple of dimension $\dim(\fg)$:
\begin{align}
\ChirHoch^0(V_k(\fg)) &= \bC,\\
\ChirHoch^1(V_k(\fg)) &\cong \fg
\qquad\text{($\dim(\fg)$ deformations)},\\
\ChirHoch^2(V_k(\fg)) &= \bC.
\end{align}
Total dimension: $\dim(\fg) + 2$. The $\ChirHoch^1$ is
isomorphic to $\fg$ as a vector space; the $\dim(\fg)$
directions correspond to the current-algebra deformations
$J^a(z) \mapsto J^a(z) + \epsilon\,\phi^a(z)$.
\end{example}

\subsection{Three Hochschild theories: chiral, topological, categorical}

\begin{warning}[Three Hochschild theories: chiral, topological, categorical]
\label{warn:e-three-hochschild}
Three distinct Hochschild-type constructions appear in the $\En$ operadic circle: chiral, topological, and categorical. The geometry of the base space
determines which of the chiral/topological/categorical theories applies.
\begin{enumerate}[label=\textup{(\roman*)}]
\item \emph{Chiral Hochschild:}
 $\ChirHoch^*(\cA)$. Input: $\Einf$-chiral algebra on a
 curve $X$. Configuration complex: $\FM_k(X)$. Output:
 $\Etwo$ structure (chiral braces), concentrated in degrees
 $\{0,1,2\}$ (Theorem~\ref{thm:e-theorem-H}).
\item \emph{Topological Hochschild:}
 $\THH(A)$. Input: $\Eone$-algebra (associative). Configuration
 complex: $\Conf_k(\bR)$ (ordered, locally constant). Output:
 $\Etwo$ structure (classical Deligne conjecture).
\item \emph{Categorical Drinfeld centre:}
 $Z(\cA\text{-mod})$. Input: $\Eone$-monoidal category of
 $\cA$-modules. Output: $\Etwo$-braided monoidal category
 (the centre carries half-braidings).
\end{enumerate}
Chiral Hochschild lives on algebraic curves; topological
Hochschild lives on $S^1$; the categorical Drinfeld centre lives in the
representation category. The three constructions give different
answers for the same algebra, and their relationship is the
content of the $\En$ operadic circle
(Section~\ref{sec:en-circle}).
\end{warning}

\subsection{The $\SCchtop$ structure on the pair
$(\cZ^{\mathrm{der}}_{\mathrm{ch}}(\cA), \cA)$}

\begin{principle}[The bar complex is $\Eone$-chiral coassociative;
$\SCchtop$ lives on the derived centre pair]
\label{princ:e-sc-location}
The ordered bar complex
$\Barord(\cA) = T^c(s^{-1}\bar{\cA})$ is an
$\Eone$-chiral coassociative coalgebra: a dg coalgebra over
$(\chirAss)^!$ with differential from OPE residues and
deconcatenation coproduct. It does \emph{not} carry $\SCchtop$
structure.

The $\SCchtop$ structure emerges on the \emph{derived chiral
centre pair}
$(\cZ^{\mathrm{der}}_{\mathrm{ch}}(\cA), \cA)$:
\begin{itemize}
\item The derived centre
 $\cZ^{\mathrm{der}}_{\mathrm{ch}}(\cA)$ is the
 \emph{closed colour} (bulk), carrying $\Etwo$ from
 $\FM_k(\bC)$.
\item The chiral algebra $\cA$ is the \emph{open colour}
 (boundary), carrying $\Eone$ from ordered configurations.
\item The bulk acts on the boundary: the action maps
 $\alpha_{k,m}\colon
 \FM_k(\bC) \times \Eone(m) \times
 (\cZ^{\mathrm{der}})^{\otimes k} \times \cA^{\otimes m}
 \to \cA$ are the mixed-sector operations.
\item No boundary-to-bulk maps exist:
 $\SCchtop(\ldots,\mathsf{o},\ldots;\,\mathsf{c})
 = \varnothing$.
\end{itemize}
The passage from bar complex to derived centre is the
chiral Hochschild construction: the convolution
$\Hom(\Barord(\cA), \cA)$ is the computational model for
$C^\bullet_{\mathrm{ch}}(\cA, \cA)$. The bar complex is the
$\Eone$ engine; the derived centre is the $\SCchtop$ output.
\end{principle}

\begin{example}[Heisenberg: $\SCchtop$ structure]
\label{ex:e-heis-sc}
For $H_k$, the derived centre pair
$(\cZ^{\mathrm{der}}_{\mathrm{ch}}(H_k), H_k)$ has closed
colour $\cZ^{\mathrm{der}} \simeq \bC^3$
(Example~\ref{ex:e-heis-chirhoch}) and open colour
$H_k$. The bulk-to-boundary action is determined by the
identity map (degree $0$), the level deformation (degree $1$),
and the obstruction class (degree $2$). All higher mixed-sector
operations vanish: the Heisenberg is SC-formal.
\end{example}

% ================================================================
% SECTION 4: The Swiss-cheese operad SC^{ch,top}
% ================================================================

\section{The Swiss-cheese operad $\SCchtop$}
\label{sec:sc-operad}

\subsection{Definition}

\begin{definition}[Holomorphic--topological Swiss-cheese operad]
\label{def:e-SC}
The two-coloured topological operad $\SCchtop$ has colours
$\{\mathsf{ch}, \mathsf{top}\}$ (closed/holomorphic and
open/topological) with three sectors:
\begin{enumerate}[label=\textup{(\roman*)}]
\item \emph{Closed sector} (all inputs and output
 closed/holomorphic):
 \[
 \SCchtop((\mathsf{ch},\ldots,\mathsf{ch});\,\mathsf{ch})
 \;=\; \FM_k(\bC).
 \]
 This is the $\Etwo$-operad parametrising holomorphic
 configurations on the complex plane.
\item \emph{Open sector} (all inputs and output
 open/topological):
 \[
 \SCchtop((\mathsf{top},\ldots,\mathsf{top});\,\mathsf{top})
 \;=\; \Eone(m).
 \]
 This is the $\Eone$-operad of ordered configurations on the
 real line.
\item \emph{Mixed sector} ($k$ closed and $m$ open inputs,
 open output):
 \[
 \SCchtop((\mathsf{ch}^k, \mathsf{top}^m);\,\mathsf{top})
 \;=\; \FM_k(\bC) \times \Eone(m).
 \]
 No mixed operations with closed output: the closed colour
 receives no information from the open colour.
\end{enumerate}
Composition is componentwise: FM insertion in $\bC$,
interval insertion in $\Eone$. The no-open-to-closed rule
is the directionality constraint
$\SCchtop(\ldots,\mathsf{top},\ldots;\,\mathsf{ch})
= \varnothing$.
\end{definition}

\begin{remark}[$\SCchtop$ is not $\Ethree$]
\label{rem:e-sc-not-e3}
The Swiss-cheese operad is two-coloured with a strict
directionality constraint. Dunn additivity
($\En \otimes E_m \simeq E_{n+m}$) applies to
\emph{single-coloured} operads. The product structure
$\FM_k(\bC) \times \Eone(m)$ in the mixed sector does
\emph{not} imply that $\SCchtop$ is $\Etwo \otimes \Eone = \Ethree$:
the directionality constraint (no open-to-closed) breaks the
symmetry required for Dunn's theorem. The passage from
$\SCchtop$ to $\Ethree$ requires additional data
(Section~\ref{sec:topologisation}).
\end{remark}

\subsection{Koszul duality of $\SCchtop$}

\begin{theorem}[Koszulity of $\SCchtop$
\textup{\cite{Livernet06}}]
\label{thm:e-sc-koszul}
The operad $\SCchtop$ is Koszul.
\end{theorem}

\begin{proposition}[Koszul dual cooperad: three sectors]
\label{prop:e-sc-dual}
The Koszul dual cooperad $(\SCchtop)^!$ has three sectors:
\begin{enumerate}[label=\textup{(\roman*)}]
\item \emph{Closed:}
 $\dim(\SCchtop)^!(n;\mathsf{c}) = (n{-}1)!$
 \quad\textup{(}$\mathrm{Lie}$ cooperad, since
 $\mathrm{Com}^! = \mathrm{Lie}$\textup{)}.
\item \emph{Open:}
 $\dim(\SCchtop)^!(m;\mathsf{o}) = m!$
 \quad\textup{(}$\mathrm{Ass}$ cooperad,
 self-Koszul-dual\textup{)}.
\item \emph{Mixed} ($k$ closed, $m$ open inputs):
 $\dim(\SCchtop)^!(k,m;\mathsf{o}) =
 (k{-}1)!\,\binom{k+m}{m}$.
\end{enumerate}
\end{proposition}

\begin{proof}
The closed sector is $\mathrm{Com}^! = \mathrm{Lie}$ with
$\dim\mathrm{Lie}(n) = (n{-}1)!$. The open sector is
$\mathrm{Ass}^! = \mathrm{Ass}$ with $\dim = m!$. The mixed
sector combines a Lie factor $(k{-}1)!$ from the holomorphic
direction with a binomial coefficient $\binom{k+m}{m}$ from the
$(k,m)$-shuffle interleaving of $k$ closed inputs among $m$
ordered open inputs on $\FM_k(\bC) \times \Eone(m)$.
\end{proof}

\begin{corollary}[$\SCchtop$ is not Koszul self-dual]
\label{cor:e-sc-not-selfdual}
The dimensions of $\SCchtop$ and $(\SCchtop)^!$ disagree in the
closed sector:
\[
\dim\,\SCchtop(n;\mathsf{c}) = 1
\quad\text{(Com)},
\qquad
\dim(\SCchtop)^!(n;\mathsf{c}) = (n{-}1)!
\quad\text{(Lie)}.
\]
Therefore $\SCchtop \not\cong (\SCchtop)^!$: the operad is
\emph{not} Koszul self-dual.

The duality \emph{functor} on $\SCchtop$-algebras is an
involution ($(P^!)^! \cong P$ for Koszul $P$), but the
operad itself is not fixed. Koszul duality exchanges
$\mathrm{Com}$ (closed) with $\mathrm{Lie}$ (closed) while
preserving $\mathrm{Ass}$ (open).
\end{corollary}

\subsection{The five presentations of $\SCchtop$}

An $\SCchtop$-algebra pair $(\cA^{\mathrm{cl}}, A^{\mathrm{op}})$
can be presented in five equivalent ways:

\begin{proposition}[Five presentations of $\SCchtop$-algebra
structure]\label{prop:e-five-presentations}
The following structures on a pair
$(\cA^{\mathrm{cl}}, A^{\mathrm{op}})$ are equivalent for
Koszul $\SCchtop$:
\begin{enumerate}[label=\textup{(\roman*)}]
\item \emph{Operadic:}
 a map of two-coloured operads
 $\SCchtop \to \End_{(\cA^{\mathrm{cl}}, A^{\mathrm{op}})}$.
\item \emph{Koszul dual:}
 an $(\SCchtop)^!$-coalgebra structure on the pair
 $(\barB(\cA^{\mathrm{cl}}), \barB(A^{\mathrm{op}}))$.
 The closed colour is a Lie coalgebra
 $(\barB_{\mathrm{Lie}}(\cA^{\mathrm{cl}}))$; the open colour is
 an associative coalgebra $(\barB_{\mathrm{Ass}}(A^{\mathrm{op}}))$.
\item \emph{Factorisation:}
 a factorisation algebra on $\bC \times \bR^{\ge 0}$
 (the upper half-space), with the interior ($\bC \times \bR^{>0}$)
 carrying the closed colour and the boundary
 ($\bC \times \{0\}$) carrying the open colour.
\item \emph{BV/BRST:}
 a BV algebra whose classical master equation encodes the
 $\SCchtop$ structure, with the antibracket splitting into
 holomorphic (closed) and topological (open) components.
\item \emph{Convolution:}
 a Maurer--Cartan element in the $\SCchtop$-convolution algebra
 $\fg^{\SCchtop}_\cA
 = \fg^{\mathrm{mod}}_\cA \times \fg^{\bR}_\cA$,
 where $\fg^{\mathrm{mod}}$ controls the closed direction
 (shadows $\kappa$, $\mathfrak{C}$, $\mathfrak{Q}$) and
 $\fg^{\bR}$ controls the open direction ($r$-matrix, KZ
 associator, Yangian deformation).
\end{enumerate}
\end{proposition}

\begin{remark}[Pentagon of equivalences]
\label{rem:e-pentagon}
The five presentations are pairwise equivalent, but all
equivalences route through the operadic hub (presentation~(i)).
No single pentagon theorem closes the circle of all five
equivalences directly. Three links have independent mathematical
content: Koszul--BV (the BV formalism realises the Koszul dual
cooperad), BV--Convolution (the antibracket encodes the
convolution Lie bracket), and Factorisation--Convolution
(the global sections functor converts a factorisation algebra
to a convolution algebra).
\end{remark}

\subsection{The Koszul dual algebra $\cA^!$ as an
$(\SCchtop)^!$-algebra}

\begin{proposition}[Koszul dual carries $(\SCchtop)^! =
(\mathrm{Lie}, \mathrm{Ass})$ structure]
\label{prop:e-koszul-dual-sc}
Let $\cA$ be a Koszul chiral algebra. The Koszul dual
$\cA^! = ((\cA^i)^\vee)$ carries the structure of an
$(\SCchtop)^!$-algebra:
\begin{enumerate}[label=\textup{(\roman*)}]
\item \emph{Closed colour:}
 the Sklyanin bracket $\{-,-\}$ makes
 $\cA^!$ a Lie algebra (from $\mathrm{Com}^! = \mathrm{Lie}$).
\item \emph{Open colour:}
 the Yangian product makes $\cA^!$ an associative algebra
 (from $\mathrm{Ass}^! = \mathrm{Ass}$).
\end{enumerate}
This is an $(\SCchtop)^!$-algebra $= (\mathrm{Lie},
\mathrm{Ass})$-algebra, \emph{not} an $\SCchtop$-algebra
$= (\mathrm{Com}, \mathrm{Ass})$-algebra.
\end{proposition}

\begin{example}[Heisenberg: $H_k^!$]
\label{ex:e-heis-koszul-dual}
The Koszul dual $H_k^! = \Sym^{\mathrm{ch}}(V^*)$ has
$\kappa(H_k^!) = -k$, so the Koszul complementarity gives
$\kappa(H_k) + \kappa(H_k^!) = k + (-k) = 0$. The Lie bracket
(Sklyanin bracket) on $H_k^!$ is trivial (abelian), reflecting
the class~G status. The Ass structure is the symmetric algebra
product, which is commutative and thus carries no ordered data
beyond the symmetric quotient.
\end{example}

% ================================================================
% SECTION 5: The topologisation theorem
% ================================================================

\section{The topologisation theorem}
\label{sec:topologisation}

\subsection{The problem}

The chiral centre theorem
(Theorem~\ref{thm:e-center-hochschild}) produces $\Etwo$
structure on the derived centre
$\cZ^{\mathrm{der}}_{\mathrm{ch}}(\cA)$, with the $\Etwo$
coming from holomorphic configurations $\FM_k(\bC)$. This
$\Etwo$ is \emph{holomorphic}: it depends on the complex
structure of the curve. The dimensional ladder places $\Ethree$
at the level of $3$-manifolds (Chern--Simons). How does the
holomorphic $\Etwo$ centre upgrade to topological $\Ethree$?

The answer requires extra structure: an inner conformal vector.
The conformal vector makes holomorphic translations $Q$-exact
in the BRST complex, killing the dependence on the complex
structure at the cohomological level. The BRST cohomology of
the factorisation algebra on $\bC$ becomes locally constant;
a locally constant factorisation algebra on $\bR^2$ is an
$\Etwo^{\mathrm{top}}$-algebra. Combined with the $\Eone$ from
the topological direction, Dunn additivity gives
$\Etwo^{\mathrm{top}} \otimes \Eone^{\mathrm{top}} =
\Ethree^{\mathrm{top}}$.

\subsection{The conformal vector}

\begin{definition}[Inner conformal vector]
\label{def:e-conformal-vector}
Let $\cA$ be a chiral algebra and $Q$ a BRST differential on the
observables of a holomorphic--topological theory on
$\bC \times \bR$ whose boundary chiral algebra is $\cA$. An
\emph{inner conformal vector} for $(\cA, Q)$ is a conformal
vector $T(z) \in \cA$ that is $Q$-exact in the algebra of bulk
observables: there exists $G(z)$ (the antighost contraction) such
that
\begin{equation}\label{eq:e-sugawara-exact}
T(z) \;=\; [Q, G(z)]
\end{equation}
in BRST cohomology. The consequence: holomorphic translations
$\partial_z \cO = [T, \cO] = [Q, [G, \cO]]$ are $Q$-exact for
every bulk observable $\cO$.
\end{definition}

\begin{example}[Affine Kac--Moody: the Sugawara element]
\label{ex:e-sugawara}
For $V_k(\fg)$ at non-critical level $k \neq -h^\vee$, the
Sugawara element
\begin{equation}\label{eq:e-sugawara}
T_{\mathrm{Sug}}
\;=\; \frac{1}{2(k + h^\vee)}\sum_a {:}J^a J_a{:}
\end{equation}
provides the inner conformal vector. The central charge is
$c = k\dim(\fg)/(k + h^\vee)$, and the antighost contraction is
\begin{equation}\label{eq:e-antighost}
G(z) \;=\; \frac{1}{2(k + h^\vee)}\sum_a
{:}J^a(z)\,\bar{c}_a(z){:},
\end{equation}
where $\bar{c}_a$ is the antighost field of the BRST complex.
\end{example}

\begin{example}[Heisenberg: the Sugawara element]
\label{ex:e-heis-sugawara}
For $H_k$ at $k \neq 0$, the Sugawara element is
$T = (1/2k)\,{:}JJ{:}$ with central charge $c = 1$. At
$k = 0$, the Sugawara element is undefined (the denominator
$2k$ vanishes) and the algebra has no conformal vector:
topologisation fails. The critical level for the Heisenberg is
$k = 0$ (there is no dual Coxeter number for $\mathrm{U}(1)$).
\end{example}

\begin{remark}[Critical level: topologisation fails]
\label{rem:e-critical-level}
At the critical level $k = -h^\vee$ for affine Kac--Moody,
the Sugawara denominator $2(k + h^\vee)$ vanishes: the Sugawara
element is undefined, $\partial_z$ is not $Q$-exact, and the
factorisation algebra retains genuine holomorphic dependence.
The $\Etwo$ structure remains holomorphic, not topological,
and the upgrade to $\Ethree$ fails. The critical level is the
precise obstruction to topologisation.
\end{remark}

\subsection{The topologisation theorem}

\begin{theorem}[Cohomological topologisation for affine Kac--Moody]
\label{thm:e-topologisation}
Let $\fg$ be a finite-dimensional simple Lie algebra and
$V_k(\fg)$ the universal affine vertex algebra at non-critical
level $k \neq -h^\vee$. The Sugawara element
$T_{\mathrm{Sug}} = (1/2(k + h^\vee))\sum_a {:}J^a J_a{:}$
provides an inner conformal vector. The derived chiral centre
$\cZ^{\mathrm{der}}_{\mathrm{ch}}(V_k(\fg))$, which carries
$\Etwo$ from the closed sector of $\SCchtop$
\textup{(}Theorem~\ref{thm:e-center-hochschild}\textup{)}, has
the following topologisation package:
\begin{enumerate}[label=\textup{(\roman*)}]
\item \textup{(Cohomological $\Ethree^{\mathrm{top}}$.)}
 The BRST cohomology
 $H_Q^\bullet(\cZ^{\mathrm{der}}_{\mathrm{ch}}(V_k(\fg)))$
 carries an $\Ethree^{\mathrm{top}}$-algebra structure:
 \begin{equation}\label{eq:e-topologisation}
 H_Q^\bullet(\cZ^{\mathrm{der}}_{\mathrm{ch}}(V_k(\fg)))
 \;\text{ is an }\;
 \Etwo^{\mathrm{top}} \otimes \Eone^{\mathrm{top}}
 = \Ethree^{\mathrm{top}}\text{-algebra.}
 \end{equation}

\item \textup{(Chain-level model.)}
 The cohomology complex
 $M_k = H_Q^\bullet(\cZ^{\mathrm{der}}_{\mathrm{ch}})$,
 with zero differential and the $\Ethree^{\mathrm{top}}$
 operations of part~\textup{(i)}, is a chain-level
 $\Ethree^{\mathrm{top}}$-algebra quasi-isomorphic to the
 derived centre. The homotopy transfer theorem controls the
 dependence on the chosen retract: transferred
 $(E_3)_\infty$-models are unique up to
 $\infty$-quasi-isomorphism.

\item \textup{(Original-complex lift: conditional.)}
 If there exists a family $G = (G_1, G_2, \ldots)$ with
 $G_1$ the Sugawara antighost contraction
 \eqref{eq:e-antighost} and $[m, G] = \partial_z$ in the
 brace deformation complex, then the original complex
 $\cZ^{\mathrm{der}}_{\mathrm{ch}}(V_k(\fg))$ itself
 admits a chain-level $\Ethree^{\mathrm{top}}$-structure.
\end{enumerate}
\end{theorem}

\begin{proof}
\emph{Part (i).}
The inner conformal vector identifies $T_{\mathrm{Sug}} =
[Q, G]$ in BRST cohomology. For any bulk observable $\cO$,
$\partial_z\cO = [T_{\mathrm{Sug}}, \cO] = [Q, [G, \cO]]$,
so $\partial_z$ is $Q$-exact. The BRST cohomology of the
factorisation algebra on $\bC$ is therefore locally constant:
it does not detect the complex structure. By the recognition
theorem of Lurie \cite{HA}, a locally constant factorisation
algebra on $\bR^2$ is an $\Etwo^{\mathrm{top}}$-algebra.
The $\Eone^{\mathrm{top}}$ from the $\bR$-direction (the open
colour of $\SCchtop$) combines with the
$\Etwo^{\mathrm{top}}$ via Dunn additivity to give
$\Ethree^{\mathrm{top}}$.

\emph{Part (ii).}
Over $\bC$, choose a splitting of cycles and boundaries in
$\cZ^{\mathrm{der}}_{\mathrm{ch}}$; this gives a strong
deformation retract onto the cohomology complex $M_k$. Part~(i)
equips $M_k$ with $\Ethree^{\mathrm{top}}$. The homotopy
transfer theorem for operadic algebras ensures that transferred
$(E_3)_\infty$-models are unique up to
$\infty$-quasi-isomorphism.

\emph{Part (iii).}
The condition $[m, G] = \partial_z$ asserts that holomorphic
translations are null-homotopic in the brace deformation complex
by an $\Ainf$-coherent family of higher homotopies. Under this
condition, the locally-constant recognition applies to the
original cochain complex (not just to cohomology), and Dunn
additivity produces $\Ethree^{\mathrm{top}}$ at chain level.

\emph{Critical level.}
At $k = -h^\vee$, the Sugawara denominator vanishes: the inner
conformal vector is undefined, $\partial_z$ is not $Q$-exact,
and the factorisation algebra retains holomorphic dependence.
The $\Etwo$ remains holomorphic, and the upgrade to $\Ethree$
fails.
\end{proof}

\begin{remark}[The conformal vector is necessary]
\label{rem:e-conformal-necessary}
Without a conformal vector, the $\SCchtop$-algebra produces a
centre with holomorphic $\Etwo$ that genuinely detects the
complex structure. The $\SCchtop$ structure is the
\emph{intermediary} between $\Eone$-chiral and $\Ethree$; the
conformal vector is the bridge. The $\Ethree$ is
$\Ethree^{\mathrm{top}}$ (topological), not
$\Ethree$-chiral: the conformal vector kills the chiral
direction.
\end{remark}

\begin{remark}[Scope of topologisation]
\label{rem:e-topologisation-scope}
Theorem~\ref{thm:e-topologisation} is proved for affine
Kac--Moody $V_k(\fg)$ at non-critical level, where the Sugawara
element provides an explicit inner conformal vector and the $3$d
holomorphic--topological theory is Chern--Simons.

For general chiral algebras with conformal vector (Virasoro,
$\cW$-algebras), the same pattern is expected but remains
conjectural:
\end{remark}

\begin{conjecture}[General topologisation]
\label{conj:e-topologisation-general}
Let $\cA$ be a chiral algebra with conformal vector $T(z)$ of
central charge $c$. If the associated holomorphic--topological
theory on $\bC \times \bR$ has a BRST differential $Q$ such
that $T = [Q, G]$ for some antighost contraction $G$, then
$H_Q^\bullet(\cZ^{\mathrm{der}}_{\mathrm{ch}}(\cA))$
carries an $\Ethree^{\mathrm{top}}$-algebra structure.

For Virasoro $\mathrm{Vir}_c$: the holomorphic--topological
theory is $3$d gravity; the BRST complex requires constructing
the gravitational path integral, which is a genuine difficulty.

For $\cW_N$: the holomorphic--topological theory is $3$d
higher-spin gravity; the BRST complex requires the full
non-linear $\cW_N$ gauge theory.
\end{conjecture}

\subsection{The Heisenberg: topologisation}

\begin{example}[Heisenberg topologisation]
\label{ex:e-heis-topologisation}
For $H_k$ at $k \neq 0$: the Sugawara element
$T = (1/2k)\,{:}JJ{:}$ provides the inner conformal vector.
The $3$d holomorphic--topological theory is abelian Chern--Simons
on $\bC \times \bR$. All three parts of
Theorem~\ref{thm:e-topologisation} hold unconditionally:
\begin{enumerate}[label=\textup{(\roman*)}]
\item BRST cohomology is $\Ethree^{\mathrm{top}}$
 (the abelian case is trivial: all operations are determined
 by the commutative product and the Lie bracket).
\item The chain-level model is the $3$-dimensional cochain
 complex $\bC \oplus \bC[-1] \oplus \bC[-2]$ with trivial
 differential and $\Ethree$ structure from the abelian
 Chern--Simons.
\item The original-complex lift is unconditional because the
 $\Ainf$ coherences are trivial (class~G: all higher operations
 vanish).
\end{enumerate}
At $k = 0$: the Sugawara element is undefined, and the algebra
is stuck at $\SCchtop$. The derived centre
$\cZ^{\mathrm{der}}_{\mathrm{ch}}(H_0)$ carries holomorphic
$\Etwo$ but not $\Ethree^{\mathrm{top}}$.
\end{example}

\subsection{Classification by topologisation status}

\begin{proposition}[Topologisation status of standard families]
\label{prop:e-topologisation-status}
\begin{center}
\renewcommand{\arraystretch}{1.15}
\begin{tabular}{llll}
\textbf{Family} & \textbf{Conformal vector} &
 \textbf{Topologisation} & \textbf{Status} \\
\hline
$H_k$, $k \neq 0$ & $T = (1/2k)\,{:}JJ{:}$ &
 $\Ethree^{\mathrm{top}}$ & proved \\
$V_k(\fg)$, $k \neq -h^\vee$ &
 $T_{\mathrm{Sug}}$ & $\Ethree^{\mathrm{top}}$ & proved \\
$V_{-h^\vee}(\fg)$ (critical) & none &
 $\SCchtop$ & stuck \\
$\mathrm{Vir}_c$ & $T(z)$ &
 $\Ethree^{\mathrm{top}}$ & conjectural \\
$\cW_N$ & $T(z)$ &
 $\Ethree^{\mathrm{top}}$ & conjectural \\
$H_0$ & none & $\SCchtop$ & stuck \\
Lattice VA, $k \neq 0$ & $T = (1/2k)\,{:}JJ{:}$ &
 $\Ethree^{\mathrm{top}}$ & proved (abelian CS)
\end{tabular}
\end{center}
\end{proposition}

% ================================================================
% SECTION 6: The E_n operadic circle
% ================================================================

\section{The $\En$ operadic circle}
\label{sec:en-circle}

\subsection{The circle}

The five theorems of modular Koszul duality live in the
symmetric bar $\BarSig(\cA)$: they are invariants surviving
the averaging map
$\mathrm{av}\colon \Barord(\cA) \to \BarSig(\cA)$.
The $\En$ operadic circle explains why the averaging map
exists, what it loses, and how the lost information can
in principle be recovered. The circle is a sequence of five
functorial operations, each shifting the operadic level by
one:
\begin{equation}\label{eq:e-circle}
\Ethree^{\mathrm{top}}(\text{bulk})
\;\xrightarrow{\;(1)\;\text{restrict}\;}
\Etwo(\text{boundary chiral})
\;\xrightarrow{\;(2)\;\Barord\;}
\Eone(\text{bar/QG})
\;\xrightarrow{\;(3)\;Z(-)\;}
\Etwo(\text{Drinfeld centre})
\;\xrightarrow{\;(4)\;\mathrm{HH}^*\;}
\Ethree^{\mathrm{top}}(\text{derived centre}).
\end{equation}

\begin{definition}[Five arrows of the operadic circle]
\label{def:e-five-arrows}
The five arrows are:
\begin{enumerate}[label=\textup{(\arabic*)}]
\item \emph{Restriction to a codimension-$2$ defect.}
 A bulk $\Ethree^{\mathrm{top}}$-algebra of observables on
 a $3$-manifold $M^3$ restricts to a factorisation algebra
 on a codimension-$2$ defect curve $C \subset M^3$. The
 normal bundle $N_{C/M}$ of rank $2$ supplies two
 directions: one holomorphic (along the curve) and one
 topological (transverse). The restricted algebra is
 $\Etwo$-chiral on $C$.
\item \emph{Ordered bar complex.}
 The $\Etwo$-chiral algebra on $C$ passes to its ordered
 bar complex $\Barord(\cA) = T^c(s^{-1}\bar{\cA})$:
 an $\Eone$-chiral coassociative coalgebra with
 deconcatenation coproduct, $R$-matrix, and (on the
 Koszul locus) Yangian structure.
\item \emph{Drinfeld centre.}
 The functor
 $Z \colon \Eone\text{-}\mathsf{Cat} \to
 \Etwo\text{-}\mathsf{Cat}$ recovers an
 $\Etwo$-braided monoidal category from $\Eone$-module
 data. This is the categorified analogue of the averaging
 map. The Drinfeld centre is the \emph{sole source of
 nontrivial braiding}: direct restriction from
 $\Ethree^{\mathrm{top}}$ to $\Etwo$ gives only
 symmetric braiding, because
 $\pi_1(\Conf_2(\bR^3))$ is trivial.
\item \emph{Higher Deligne conjecture.}
 The topological Hochschild cohomology of an $\Etwo$-algebra carries
 $\Ethree^{\mathrm{top}}$-structure. This is the higher
 Deligne conjecture, proved by
 Francis~\cite{Francis2013}: $\HH^*(A)$ for an
 $\En$-algebra $A$ carries an $E_{n+1}$-structure.
\item \emph{Closing.}
 The fifth arrow closes the circle: the
 $\Ethree^{\mathrm{top}}$-algebra obtained by traversing
 arrows $(1)$--$(4)$ should be equivalent to the original
 bulk $\Ethree^{\mathrm{top}}$-algebra. This is the
 content of the closing conjecture
 (Section~\ref{sec:closing-conjecture}).
\end{enumerate}
\end{definition}

\begin{theorem}[Arrow 3 as the sole source of braiding]
\label{thm:e-arrow3-braiding}
Let $A$ be an $\Ethree^{\mathrm{top}}$-algebra.
\begin{enumerate}[label=\textup{(\roman*)}]
\item Restriction of $A$ to a codimension-$2$ locus gives
 an $\Etwo$-algebra $A|_C$ whose braiding is
 \emph{symmetric}: the monodromy representation factors
 through $\pi_1(\Conf_2(\bR^3)) = \{1\}$, the trivial
 group.
\item The Drinfeld centre $Z(\cA\textup{-mod})$ of the
 $\Eone$-monoidal category of $\cA$-modules carries a
 half-braiding that is generically \emph{non-symmetric}:
 the braiding comes from the $R$-matrix of the ordered
 bar complex, which records the nontrivial monodromy of
 $\pi_1(\Conf_2(\bC)) = \bZ$ (the braid group on $2$
 strands).
\item The non-symmetric braiding in the Drinfeld centre is
 therefore not inherited from the original $\Ethree$
 restriction; it is created by arrow~$3$. For
 $\Einf$-chiral algebras, the braiding is symmetric
 \textup{(}the $R$-matrix is trivial\textup{)}. For
 $\Eone$-chiral algebras, the braiding is non-symmetric
 and encodes the KZ associator
 $\Phi_{\mathrm{KZ}} \in \exp(\hat{\mathfrak{t}}_3)$.
\end{enumerate}
\end{theorem}

\begin{proof}
Part~(i): the fundamental group $\pi_1(\Conf_2(\bR^3))$ is
trivial because $\bR^3 \setminus \{0\}$ is simply connected
(the linking sphere $S^2$ carries no $1$-cycles). The
representation of the braid group on the $\Etwo$-restricted
modules therefore factors through $\Sigma_2$: the braiding
is a symmetric involution.

Part~(ii): the Drinfeld centre $Z(\cA\text{-mod})$ consists
of pairs $(M, \gamma)$ where $\gamma_N \colon M \otimes N
\xrightarrow{\sim} N \otimes M$ is a half-braiding natural
in $N$. When $\cA$ is the module category of the ordered
bar complex, the half-braiding is conjugation by the
$R$-matrix: $\gamma_N = R_{M,N} \circ \tau$. For
$\Eone$-chiral algebras with nontrivial $R$-matrix, the
braiding is non-symmetric.

Part~(iii): since $\Ethree$ restriction gives only symmetric
braiding and the Drinfeld centre produces non-symmetric
braiding, the non-symmetric part is created by arrow~$3$.
This is the mathematical content of the statement that the
Drinfeld centre is the sole source of nontrivial braiding.
\end{proof}

\subsection{The $\SCchtop$ intermediary}

Between the $\Eone$-chiral world and the
$\Ethree^{\mathrm{top}}$ bulk sits the Swiss-cheese operad
$\SCchtop$ (Section~\ref{sec:sc-operad}). The passage
from $\SCchtop$ to $\Ethree^{\mathrm{top}}$ requires an
inner conformal vector
(Section~\ref{sec:topologisation}). Without such a vector,
the $\SCchtop$-algebra cannot topologise: the two colours
(holomorphic and topological) remain distinct, and Dunn
additivity does not apply.

\begin{proposition}[The $\SCchtop$ intermediary in the circle]
\label{prop:e-sc-intermediary}
The $\En$ operadic circle factors through $\SCchtop$:
\[
\Ethree^{\mathrm{top}}
\;\xrightarrow{\;(1)\;}
\Etwo^{\mathrm{ch}}
\;\xrightarrow{\;(2)\;}
\Eone^{\mathrm{ch}}
\;\xrightarrow{\;\SCchtop\;}
(\Etwo^{\mathrm{ch}},\,\Eone^{\mathrm{top}})
\;\xrightarrow{\;\nu\;}
\Ethree^{\mathrm{top}}.
\]
The penultimate stage is the $\SCchtop$-pair
$(\cZ^{\mathrm{der}}_{\mathrm{ch}}(\cA), \cA)$.
The last arrow requires the inner conformal vector $\nu$;
when $\nu$ exists, the two SC colours collapse to a single
$\Ethree^{\mathrm{top}}$ structure.
\end{proposition}

\begin{proof}
The Drinfeld centre (arrow~$3$) and the higher Deligne
conjecture (arrow~$4$) combine to produce the $\SCchtop$-pair:
the derived centre $\cZ^{\mathrm{der}}_{\mathrm{ch}}(\cA)$
carries the closed $\Etwo$ from the FM configuration spaces,
and $\cA$ carries the open $\Eone$ from ordered configurations.
The topologisation theorem
(Theorem~\ref{thm:e-topologisation}) provides the final
arrow.
\end{proof}

\subsection{The formality bridge and its failure}

\begin{theorem}[Formality bridge for $\Einf$-chiral algebras]
\label{thm:e-formality-bridge}
For an $\Einf$-chiral algebra $\cA$ \textup{(}the entire
standard families: Heisenberg, affine Kac--Moody, Virasoro,
$\cW_N$, and all free-field algebras\textup{)}, the averaging
map is a quasi-isomorphism:
\[
\mathrm{av}\colon
\Barord(\cA)
\;\xrightarrow{\;\sim\;}
\BarSig(\cA).
\]
The $R$-matrix is trivial \textup{(}the braiding on the
ordered bar is symmetric\textup{)}, the Yangian is
cocommutative, and the averaging map loses nothing.
\end{theorem}

\begin{proof}
The $\Einf$-chiral condition means the OPE is
$\Sigma_k$-equivariant: the factorisation $\cD$-module
descends to the symmetric Ran space. The $R$-matrix
$S(z) = \id$ is the identity, and the $\Sigma_k$-action
on the ordered bar complex is free. The averaging map
takes $\Sigma_k$-coinvariants of a free action, which is
a quasi-isomorphism by the standard acyclicity of the
group algebra.
\end{proof}

\begin{theorem}[Formality failure for $\Eone$-chiral algebras]
\label{thm:e-formality-failure}
For a genuinely $\Eone$-chiral algebra \textup{(}Yangian,
Etingof--Kazhdan quantum vertex algebra\textup{)}, the
averaging map is \emph{not} a quasi-isomorphism:
\[
\dim H^n(\Barord(\cA))
\;=\; n!
\quad\text{at degree $n$},
\qquad
\dim H^n(\BarSig(\cA))
\;=\; 1.
\]
The ordered bar carries Arnold braid-group representations
(from $\pi_1(\Conf_k(\bC)) = B_k$) that the symmetric bar
annihilates.
\end{theorem}

\begin{proof}
The $\Eone$-chiral condition means the $R$-matrix
$S(z) \neq \id$ is a nontrivial solution of the quantum
Yang--Baxter equation. The $\Sigma_k$-action on the
ordered bar complex is twisted by $S(z)$, and the kernel
of $\mathrm{av}$ records the braid-group representations.
The dimension $n!$ at degree $n$ is the dimension of the
regular representation of $\Sigma_n$; the coinvariant
quotient collapses to dimension $1$.
\end{proof}

\begin{remark}[What the five theorems see and what they miss]
\label{rem:e-five-thms-see}
The five theorems of modular Koszul duality (A, B, C, D, H)
extract the $\Sigma_n$-invariant content of the ordered bar:
the scalar $\kappa$, given by $\mathrm{av}(r(z))$ in abelian and
scalar families and by $\mathrm{av}(r(z))+\dim(\fg)/2$ for
non-abelian affine Kac--Moody, the genus
anomaly $\mathrm{obs}_g = \kappa \cdot \lambda_g$
(uniform-weight), and the shadow tower. For $\Einf$-chiral
algebras this is the complete picture. For $\Eone$-chiral
algebras, the five theorems see only the shadow, while the
ordered bar carries the $R$-matrix, the KZ associator, and
the full braid-group representation.
\end{remark}

% ================================================================
% SECTION 7: The E_3 identification theorem
% ================================================================

\section{The $\Ethree$ identification theorem}
\label{sec:e3-identification}

\subsection{The problem}

The topologisation theorem
(Theorem~\ref{thm:e-topologisation}) produces an
$\Ethree^{\mathrm{top}}$-algebra on the BRST cohomology of
the derived chiral centre
$\cZ^{\mathrm{der}}_{\mathrm{ch}}(V_k(\fg))$ for simple
$\fg$ at non-critical level. This is a formal deformation of
the Chevalley--Eilenberg cochains
$C^*(\fg) = \Sym(\fg^\vee[-1])$, parametrised by
$\lambda = k + h^\vee$. There is a second source of
$\Ethree^{\mathrm{top}}$-algebras: the perturbative
quantisation of Chern--Simons theory on $\bR^3$ with gauge
algebra $\fg$, carried out by
Costello--Francis--Gwilliam~\cite{CFG25}. Both are
deformations of the same classical limit. Do they coincide?

\subsection{The $\Pthree$ reduction}

\begin{theorem}[$\Ethree$ formality
\textup{\cite{Kon99,Tamarkin03,FresseWillwacher20}}]
\label{thm:e-e3-formality}
The operad $\Ethree$ is formal over $\bQ$: there is a
zig-zag of quasi-isomorphisms of operads
\[
\Ethree
\;\xleftarrow{\sim}\; \cdots \;\xrightarrow{\sim}\;
H_*(\Ethree) \;=\; \Pthree.
\]
An $\Ethree$-algebra structure on a cochain complex is
therefore determined, up to homotopy, by the induced
$\Pthree$-algebra structure on cohomology.
\end{theorem}

\begin{remark}[From $\Ethree$ to $\Pthree$]
\label{rem:e-e3-to-p3}
The $\Pthree$-algebra structure on
$H^*(\cZ^{\mathrm{der}}_{\mathrm{ch}}(V_k(\fg)))
\cong C^*(\fg)$ consists of:
\begin{enumerate}[label=\textup{(\roman*)}]
\item a graded-commutative product (the cup product,
 $\Sym(\fg^\vee[-1])$);
\item a degree-$(1{-}3) = -2$ Lie bracket
 $\{-,-\}$ (the Schouten--Nijenhuis bracket on polyvector
 fields, or equivalently the linearisation of the
 Chern--Simons interaction).
\end{enumerate}
The $\Pthree$ bracket is determined by its restriction to
generators: $\{\phi_a, \phi_b\} = \lambda\,f_{ab}{}^c\,\phi_c$,
where $\lambda = k + h^\vee$ and $f_{ab}{}^c$ are the
structure constants of $\fg$. The coefficient $\lambda$ comes
from the formal disk computation: the chiral derived centre
at level $k$ deforms the classical cochains by exactly
$\lambda = k + h^\vee$ times the Cartan $3$-form.
\end{remark}

\subsection{The identification}

\begin{theorem}[Identification of $\Ethree$-deformation families
for simple $\fg$]
\label{thm:e-e3-identification}
Let $\fg$ be a simple finite-dimensional Lie algebra,
$V_k(\fg)$ the universal affine vertex algebra at non-critical
level $k \neq -h^\vee$, and
$\lambda = k + h^\vee$. Then the derived chiral centre
$\cZ^{\mathrm{der}}_{\mathrm{ch}}(V_k(\fg))$ and the
Costello--Francis--Gwilliam $\Ethree$-algebra
$\cA^\lambda$ from perturbative Chern--Simons are isomorphic
as formal deformation families of $\Ethree$-algebras:
\begin{equation}\label{eq:e-e3-identification}
\bigl\{
\cZ^{\mathrm{der}}_{\mathrm{ch}}(V_k(\fg))
\bigr\}_{k + h^\vee \in \lambda\bC[[\lambda]]}
\;\simeq\;
\bigl\{
\cA^\lambda
\bigr\}_{\lambda \in \lambda H^3(\fg)[[\lambda]]}.
\end{equation}
\end{theorem}

\begin{proof}
The proof proceeds in four steps.

\emph{Step 1: $\Ethree$ formality reduces to $\Pthree$.}
By Theorem~\ref{thm:e-e3-formality}, both sides are
determined by their $\Pthree$-structures on the common
cohomology $C^*(\fg) = \Sym(\fg^\vee[-1])$. A formal
deformation of a $\Pthree$-algebra with cohomology
$\Sym(V[-1])$ is classified, order by order in $\lambda$,
by the cohomology group
$H^*(\Sym(V[-1]),\, \Sym(V[-1]))$ in the $\Pthree$-deformation
complex. For the Lie algebra cochains, this deformation complex
is controlled by $H^3(\fg)$ at leading order.

\emph{Step 2: one-dimensionality of $H^3(\fg)$.}
For $\fg$ simple, $H^3(\fg;\bC)$ is one-dimensional, generated
by the Cartan $3$-form
$\omega_3(X,Y,Z) = \langle X, [Y,Z]\rangle$, where
$\langle\cdot,\cdot\rangle$ is the Killing form. The
deformation base is therefore one-dimensional at each order
in $\lambda$: any $\Pthree$-deformation of
$C^*(\fg)$ is determined by a single scalar at each order.

\emph{Step 3: formal disk comparison.}
Both deformations restrict to the formal disk
$D = \Spec\bC[[z]]$. On $D$, the chiral derived centre is
computed from the bar complex $\barB_D(V_k(\fg))$
(Definition~\ref{def:e-bar-disk}), and the CFG algebra
is computed from the $1$-loop Feynman integral on $\bR^3$
localised near a point. Both produce the same $\Pthree$
bracket $\{\phi_a, \phi_b\} = \lambda\,f_{ab}{}^c\,\phi_c$
on generators. Since the deformation at each order is
controlled by a single scalar (Step~2), and the formal disk
determines that scalar (by evaluation at a point), the two
families coincide order by order.

\emph{Step 4: global extension.}
The local isomorphism (Step~3) extends to a global isomorphism
of filtered $\Ethree$-algebras because both deformations are
defined over the same one-dimensional base
$\lambda H^3(\fg)[[\lambda]]$ and agree at each order. The
filtration is by powers of $\lambda$; the isomorphism is the
identity at each graded piece.
\end{proof}

\begin{remark}[The role of $\dim H^3(\fg) = 1$]
\label{rem:e-h3-role}
The proof uses $\dim H^3(\fg) = 1$ at Step~2 in an essential
way. For non-simple $\fg$ (e.g.\
$\fg = \mathfrak{gl}_N = \mathfrak{sl}_N \oplus
\mathfrak{gl}_1$), the deformation space is
multi-dimensional: $\dim \Sym^2(\fg^*)^\fg = r +
\binom{\dim\mathfrak{z}+1}{2}$, where $r$ is the number of
simple summands and $\mathfrak{z}$ is the centre. The formal
disk comparison must then determine all independent scalars,
which requires matching all independent invariant bilinear forms
(e.g.\ $B_{\mathrm{tr}}(X,Y) = \mathrm{tr}(XY)$ and
$B_{\mathrm{ab}}(X,Y) = \mathrm{tr}(X)\mathrm{tr}(Y)$ for
$\mathfrak{gl}_N$). The identification extends to
$\mathfrak{gl}_N$ for all $N \ge 1$ via this refined matching,
but the argument requires case analysis at each $N$.
\end{remark}

\begin{corollary}[The $\En$ circle partly closes for simple $\fg$]
\label{cor:e-circle-partly-closes}
For simple $\fg$ at non-critical level:
\begin{enumerate}[label=\textup{(\roman*)}]
\item The bulk $\Ethree^{\mathrm{top}}$-algebra
 \textup{(}Chern--Simons, arrow origin\textup{)} maps to the
 boundary $\Etwo$-chiral algebra $V_k(\fg)$
 \textup{(}arrow~$1$\textup{)}.
\item The ordered bar complex extracts the $\Eone$ data
 \textup{(}arrow~$2$\textup{)}.
\item The Drinfeld centre recovers $\Etwo$
 \textup{(}arrow~$3$\textup{)}.
\item The higher Deligne conjecture produces an
 $\Ethree^{\mathrm{top}}$-algebra
 \textup{(}arrow~$4$\textup{)}.
\item The $\Ethree$ identification
 \textup{(}Theorem~\ref{thm:e-e3-identification}\textup{)}
 establishes that this $\Ethree^{\mathrm{top}}$-algebra
 is isomorphic to the original CFG algebra, closing the
 circle at the level of formal deformation families.
\end{enumerate}
\end{corollary}

\begin{remark}[Sugawara constraint]
\label{rem:e-sugawara-constraint}
The identification~\eqref{eq:e-e3-identification} carries
an additional constraint: the Sugawara element
$T_{\mathrm{Sug}} = (1/2(k+h^\vee))\sum_a{:}J^aJ_a{:}$
must be $Q$-exact in the CFG BRST complex. This is verified
for $V_k(\fg)$ at non-critical level by the explicit
antighost contraction
(Example~\ref{ex:e-sugawara}). At the critical level
$k = -h^\vee$, the Sugawara element diverges, topologisation
fails, and the $\Ethree$ identification does not apply.
\end{remark}

\begin{remark}[Extension to $\mathfrak{gl}_N$]
\label{rem:e-gl-n-extension}
For $\fg = \mathfrak{gl}_N$, the $\Ethree$ identification
extends via two independent invariant bilinear forms
$B_{\mathrm{tr}}$ and $B_{\mathrm{ab}}$. The deformation
base is two-dimensional at each order, parametrised by
$(\lambda_1, \lambda_2) \in \bC^2$, where
$\lambda_1 = k_1 + N$ controls the $\mathfrak{sl}_N$ factor
and $\lambda_2 = k_2$ controls the $\mathfrak{gl}_1$ factor.
The formal disk comparison determines both scalars
independently: choosing $a, b \in \mathfrak{sl}_N$ with
$\mathrm{tr}(a) = \mathrm{tr}(b) = 0$ and
$\mathrm{tr}(ab) \neq 0$ extracts $\lambda_1$; choosing
$a = b = I_N/N$ extracts $\lambda_2$. The isomorphism
\[
\bigl\{
\cZ^{\mathrm{der}}_{\mathrm{ch}}(V_{k_1,k_2}(\mathfrak{gl}_N))
\bigr\}
\;\simeq\;
\bigl\{
\cA^{\lambda_1,\lambda_2}
\bigr\}
\]
holds over the two-dimensional base, order by order in the
joint $(\lambda_1,\lambda_2)$-adic filtration.
\end{remark}

% ================================================================
% SECTION 8: Five notions of E_1-chiral algebra
% ================================================================

\section{Five notions of $\Eone$-chiral algebra}
\label{sec:five-notions}

\subsection{The problem}

The phrase ``$\Eone$-chiral algebra'' admits at least five
inequivalent formalizations. Each captures a different aspect
of the same informal idea (an associative algebra object in
the chiral world), but they lead to different derived centres,
different obstruction theories, and different relationships
to the $\En$ operadic circle. This section catalogues the five
notions, states the known relationships, and identifies the
open comparisons.

\subsection{The five notions}

\begin{definition}[Five notions of $\Eone$-chiral algebra]
\label{def:e-five-notions}
Let $X$ be a smooth complex curve. An \emph{$\Eone$-chiral
algebra on $X$} can mean any of the following:
\begin{enumerate}[label=\textup{(\Alph*)}]
\item \emph{Strict $\chirAss$-algebra.}
 An algebra over the associative chiral operad
 $\chirAss$, i.e.\ a $\cD$-module $\cA$ on $X$ with a
 binary product
 $\mu \colon \cA \boxtimes \cA \to \Delta_!\cA$
 satisfying strict associativity:
 $\mu \circ (\mu \otimes \id) = \mu \circ (\id \otimes \mu)$.
 This is the $\cD$-module-theoretic formulation.

\item \emph{$\Ainf$-algebra in $\End^{\mathrm{ch}}_\cA$.}
 An operad map
 $\Ainf \to \End^{\mathrm{ch}}_\cA$, where
 $\End^{\mathrm{ch}}_\cA$ is the chiral endomorphism operad.
 The structure maps are
 $m_k^{\mathrm{ch}} \colon
 \cA^{\otimes k} \to
 \cA((\lambda_1))\cdots((\lambda_{k-1}))$
 satisfying the Stasheff identities. This is the working
 definition of the present paper and the monograph.
 Strictification ($m_k = 0$ for $k \ge 3$) recovers
 Notion~(A).

\item \emph{Etingof--Kazhdan quantum vertex algebra.}
 A pair $(\cA, S(z))$ where $S(z)$ is a vertex $R$-matrix
 satisfying the five EK axioms: $S$-locality, quantum
 Yang--Baxter equation, unitarity, shift condition, and
 hexagon. The $R$-matrix is the primary datum.

\item \emph{$\Ainf$-algebra object in $\Eone$-chiral algebras.}
 A double $\Ainf$ structure: an inner chiral $\Ainf$ from
 Notion~(B) controlling the algebra, and an outer $\Ainf$
 controlling higher coherences between $\Eone$-chiral
 algebra maps. This arises naturally when studying functors
 and natural transformations in the $\Eone$-chiral world.

\item \emph{Factorisation algebra on
 $\Ran^{\mathrm{ord}}(X)$.}
 The $\cD$-module-theoretic incarnation: a factorisation
 algebra on the ordered Ran space
 $\Ran^{\mathrm{ord}}(X) = \coprod_{k \ge 1}
 \Conf_k^{\mathrm{ord}}(X)$, where the ordered
 factorisation structure encodes the $\Eone$ data
 geometrically.
\end{enumerate}
\end{definition}

\begin{theorem}[Comparison of notions (B) and (C) on the
Koszul locus]
\label{thm:e-bc-comparison}
On the Koszul locus \textup{(}the locus where the bar
complex has cohomology concentrated in degree $1$\textup{)},
Notions \textup{(B)} and \textup{(C)} determine each other
via a Drinfeld associator
$\Phi \in \exp(\hat{\mathfrak{t}}_3)$:
\begin{enumerate}[label=\textup{(\roman*)}]
\item \textup{(B)} $\Rightarrow$ \textup{(C):}
 the $\Ainf$ structure maps $m_k^{\mathrm{ch}}$ determine
 the $R$-matrix $S(z)$ via the degree-$2$ bar cohomology
 class, and the Stasheff identities imply the five EK axioms.
\item \textup{(C)} $\Rightarrow$ \textup{(B):}
 the EK quantum vertex algebra $(\cA, S(z))$ determines the
 $\Ainf$ structure maps by inverting the bar-cobar
 adjunction: the cobar of the bar coalgebra (with $R$-twisted
 differential) recovers the $\Ainf$ structure up to
 quasi-isomorphism.
\end{enumerate}
The comparison depends on the choice of associator $\Phi$:
different choices give gauge-equivalent $\Ainf$ structures.
On BRST cohomology, the $\Ainf$ structure is
associator-independent \textup{(}the $\GRT_1$ action on the
$\Pthree$ bracket is trivial for simple $\fg$\textup{)}.
At chain level, the associator dependence is genuine.
\end{theorem}

\begin{proof}
The direction (B) $\Rightarrow$ (C) is the content of the bar
construction: the degree-$2$ class in $H^*(\Barord(\cA))$ is
the vertex $R$-matrix $S(z)$, and the Stasheff identity at
degree $3$ is the quantum Yang--Baxter equation. The remaining
EK axioms (unitarity, shift, hexagon) follow from the
augmentation, the spectral parameter, and the associativity
coherences of the $\Ainf$ structure.

The direction (C) $\Rightarrow$ (B) uses the Etingof--Kazhdan
reconstruction theorem: the quantum vertex algebra axioms
determine a pro-nilpotent completion of the bar complex, and
the cobar of this completion gives the $\Ainf$ structure.
The Drinfeld associator $\Phi$ enters in the passage from the
topological $R$-matrix (a solution of QYBE on vector spaces)
to the chiral $R$-matrix (a solution on $\cD$-modules): the
associator controls the regularisation of the formal power
series expansion.
\end{proof}

\begin{warning}[Each notion has its own derived centre]
\label{warn:e-five-derived-centres}
The five notions of $\Eone$-chiral algebra lead to five
a priori distinct derived centres:
\begin{enumerate}[label=\textup{(\Alph*)}]
\item $\cZ^{(A)} = \RHom_{(\chirAss)^{e}}(\cA, \cA)$:
 the operadic chiral Hochschild cochains of the strict
 $\chirAss$-algebra.
\item $\cZ^{(B)} = C^\bullet_{\mathrm{ch}}(\cA, \cA)$:
 the chiral Hochschild cochain complex, using the $\Ainf$
 structure.
\item $\cZ^{(C)} = \RHom_{(\cA,S)}(\cA, \cA)$:
 the derived endomorphisms of $\cA$ as an EK quantum vertex
 algebra.
\item $\cZ^{(D)}$: the homotopy coherent derived centre,
 using the double $\Ainf$ structure.
\item $\cZ^{(E)} = \int_{\Ran^{\mathrm{ord}}(X)} \cA$:
 the factorisation homology of the ordered factorisation
 algebra.
\end{enumerate}
The five derived centres are not known to coincide in general.
On the Koszul locus, \textup{(B)} and \textup{(C)} give the
same derived centre \textup{(}by
Theorem~\ref{thm:e-bc-comparison}\textup{)}. The comparison
of the remaining notions is open.
\end{warning}

\begin{remark}[Symmetric input simplifies everything]
\label{rem:e-symmetric-simplifies}
For $\Einf$-chiral algebras (all standard vertex algebras),
the five notions collapse: the $R$-matrix is trivial
($S(z) = \id$), the $\Ainf$ structure is strict ($m_k = 0$
for $k \ge 3$ on cohomology), and the factorisation algebra
descends to the unordered Ran space. All five derived centres
coincide. The distinction matters only for genuinely
$\Eone$-chiral algebras (Yangians, EK quantum vertex algebras).
\end{remark}

\begin{remark}[The monograph uses Notion~(B)]
\label{rem:e-monograph-notion-b}
Every theorem in this paper stated for ``an $\Eone$-chiral
algebra $\cA$'' means a Notion~(B) structure unless explicitly
qualified. Notion~(B) is the natural setting for the bar-cobar
adjunction: the $\Ainf$ structure maps are the components of
the bar differential, and the Stasheff identities are $d^2 = 0$.
\end{remark}

% ================================================================
% SECTION 9: The chiral quantum group equivalence
% ================================================================

\section{The chiral quantum group equivalence}
\label{sec:chiral-qg}

\subsection{The problem}

The ordered bar complex
$\Barord(\cA) = T^c(s^{-1}\bar{\cA})$
is a single object from which three structures are recovered:
the vertex $R$-matrix $S(z)$ (braiding data from the
degree-$2$ collision residue), the chiral $\Ainf$-structure
maps $m_k^{\mathrm{ch}}$ (higher associativity from boundary
strata of associahedra), and the chiral coproduct
$\Delta^{\mathrm{ch}}$ (coalgebra data from deconcatenation).
The three structures determine each other on the Koszul locus.

\subsection{The equivalence theorem}

\begin{theorem}[Chiral quantum group equivalence]
\label{thm:e-chiral-qg-equiv}
The following three structures on an $\Eone$-chiral
algebra $\cA$ on the Koszul locus determine each other,
up to the choice of a Drinfeld associator
$\Phi \in \exp(\hat{\mathfrak{t}}_3)$:
\begin{enumerate}[label=\textup{(\Roman*)}]
\item \emph{$R$-matrix datum:}
 $\Eone$-chiral algebra with vertex $R$-matrix $S(z)$
 satisfying $S$-locality, the quantum Yang--Baxter
 equation, unitarity, shift condition, and hexagon axiom
 \textup{(}the Etingof--Kazhdan axioms\textup{)}.
\item \emph{$\Ainf$ datum:}
 $\Ainf$-algebra in the chiral endomorphism operad
 $\End^{\mathrm{ch}}_\cA$, with structure maps
 $m_k^{\mathrm{ch}} \colon
 \cA^{\otimes k} \to
 \cA((\lambda_1))\cdots((\lambda_{k-1}))$
 satisfying the Stasheff identities.
\item \emph{Coproduct datum:}
 $\Eone$-chiral algebra with compatible chiral coproduct
 $\Delta^{\mathrm{ch}} \colon
 \cA \to (\cA \,\hat{\otimes}\, \cA)((z))$,
 coassociative up to conjugation by $\Phi$.
\end{enumerate}
The equivalences form a triangle:
\begin{enumerate}[label=\textup{(\roman*)}]
\item $\alpha\colon$ \textup{(I)} $\to$ \textup{(II):}
 The $R$-matrix determines the $\Ainf$ structure via the
 ordered bar complex: the degree-$k$ component of the bar
 differential, twisted by $S(z)$, gives $m_k^{\mathrm{ch}}$.
\item $\beta\colon$ \textup{(II)} $\to$ \textup{(III):}
 The $\Ainf$ structure determines the coproduct via the
 deconcatenation coalgebra structure on the bar complex:
 the coproduct $\Delta^{\mathrm{ch}}$ is the desuspended
 deconcatenation $\Delta$ of $\Barord(\cA)$, transferred
 through the bar-cobar quasi-isomorphism.
\item $\gamma\colon$ \textup{(III)} $\to$ \textup{(I):}
 The coproduct determines the $R$-matrix via the universal
 $R$-matrix of the Drinfeld double: the $R$-matrix is
 the image of the canonical element
 $\sum e_i \otimes e^i \in
 \cA \,\hat{\otimes}\, \cA^!$ under the coproduct.
\end{enumerate}
\end{theorem}

\begin{proof}
The functor $\alpha$ is the bar construction: the ordered bar
complex $\Barord(\cA)$ with $R$-twisted differential has
degree-$k$ components that are the $\Ainf$ structure maps. The
Stasheff identities are equivalent to $d^2 = 0$, which holds
by the Arnold relation (Theorem~\ref{thm:e-d2-totaro}).

The functor $\beta$ uses the coalgebra structure: the
deconcatenation coproduct on $\Barord(\cA) = T^c(s^{-1}\bar{\cA})$
is coassociative (the cofree tensor coalgebra is
coassociative). The bar-cobar quasi-isomorphism transfers
this coproduct to $\cA$, producing $\Delta^{\mathrm{ch}}$.
The transfer is exact on the Koszul locus (the bar
cohomology is concentrated in degree $1$, so higher transfer
terms vanish).

The functor $\gamma$ is the universal $R$-matrix construction:
the Drinfeld double
$U_\cA = \cA \bowtie \cA^!$ carries a canonical element
$\cR = \sum e_i \otimes e^i$, where $\{e_i\}$ and $\{e^i\}$
are dual bases of $\cA$ and $\cA^!$ respectively. The image of
$\cR$ under $\Delta^{\mathrm{ch}} \otimes \id$ gives the
vertex $R$-matrix.

The triangle closes because $\gamma \circ \beta \circ \alpha$
is the identity: starting from $S(z)$, constructing the $\Ainf$
structure, extracting the coproduct, and recovering the
$R$-matrix gives back $S(z)$.
\end{proof}

\begin{remark}[Associator dependence]
\label{rem:e-associator-dependence}
The equivalences depend on the choice of Drinfeld associator
$\Phi$. Different choices give gauge-equivalent structures:
the $R$-matrices differ by a conjugation, the $\Ainf$
structures differ by an $\Ainf$ quasi-isomorphism, and the
coproducts differ by a twist. On BRST cohomology, the
bar-side invariants ($\kappa$, shadow tower) are
associator-free; the cobar-side and coproduct-side carry
genuine associator dependence.
\end{remark}

\subsection{The $\mathfrak{gl}_N$ tower}

\begin{theorem}[$\mathfrak{gl}_N$ chiral quantum group for
all $N \ge 1$]
\label{thm:e-glN-chiral-qg}
For every $N \ge 1$, the principal $\cW$-algebra
$\cW_N$ carries a chiral quantum group datum:
\begin{enumerate}[label=\textup{(\Roman*)}]
\item \emph{$R$-matrix:} the Yang $R$-matrix
 $R(u) = u\,I + \Psi\,P$, where $I$ is the $N \times N$
 identity, $P$ is the permutation matrix, and $\Psi$ is the
 coupling parameter.
\item \emph{Coproduct:} the Drinfeld coproduct
 $\Delta_z(T(u)) = T(u) \cdot T(u - z)$ as $N \times N$
 matrix multiplication, where $T(u)$ is the generating matrix
 $T(u) = \sum_{s=0}^{N} \psi_s\,u^{-s}$ with
 $\psi_s = e_s(\Lambda_1,\ldots,\Lambda_N)$ the elementary
 symmetric polynomials in the Miura operators.
\item \emph{RTT relations:} for $N \ge 2$,
 \[
 R(u - v)\,(T(u) \otimes I)\,(I \otimes T(v))
 \;=\;
 (I \otimes T(v))\,(T(u) \otimes I)\,R(u - v),
 \]
 which is nontrivial precisely when $N \ge 2$
 \textup{(}at $N = 1$, the RTT relation is the scalar
 identity $R(u-v) = R(u-v)$\textup{)}.
\end{enumerate}
Concrete verifications run through $N = 3$.
\end{theorem}

\begin{proof}
The $R$-matrix is the Yang solution of the QYBE for
$\mathfrak{gl}_N$. The Drinfeld coproduct is an algebra
homomorphism of the Yangian $Y(\widehat{\mathfrak{gl}}_N)$:
$\Delta_z(T(u)) = T(u) \cdot T(u-z)$ expresses the
multiplicativity of the generating matrix under the
spectral-parameter shift. The RTT relation is the
compatibility of the $R$-matrix with the coproduct, proved
by direct verification at each $N$: the relation amounts to
the commutativity of the monodromy matrices of the KZ
connection at separated spectral parameters.

The spectral coassociativity uses shifted parameters:
$(\Delta_{z_1} \otimes \id) \circ \Delta_{z_1+z_2}
= (\id \otimes \Delta_{z_2}) \circ \Delta_{z_1}$.
\end{proof}

\subsection{Miura universality}

\begin{theorem}[Miura cross-term universality]
\label{thm:e-miura-universality}
For all spins $s \ge 2$, the spectral coproduct on the
spin-$s$ primary generator $W_s$ of
$\cW_{1+\infty}[\Psi]$ satisfies
\begin{equation}\label{eq:e-miura-cross}
\Delta_z(W_s)
\;=\;
W_s \otimes 1 \;+\; 1 \otimes W_s
\;+\;
\frac{\Psi - 1}{\Psi}
\bigl(J \otimes W_{s-1} + W_{s-1} \otimes J\bigr)
\;+\; (\textup{lower-spin and composite corrections})
\;+\; (\textup{spectral tail in $z$}).
\end{equation}
The coefficient $(\Psi - 1)/\Psi$ of the primary
cross-term $J \otimes W_{s-1} + W_{s-1} \otimes J$ is
independent of the spin $s$.
\end{theorem}

\begin{proof}
The proof uses the Prochazka--Rapcak Miura factorisation
$T(u) = \prod_{i=1}^{N} (u + \Lambda_i)$, where the
$\Lambda_i$ are the Miura operators. The Drinfeld coproduct
acts on $T(u)$ as matrix multiplication:
$\Delta_z(T(u)) = T(u) \cdot T(u-z)$.

The generators $\psi_s = e_s(\Lambda_1,\ldots,\Lambda_N)$
are elementary symmetric polynomials. The Drinfeld formula
at $u^{-s}$ gives
$\Delta_z(\psi_s) = \sum_{a+b=s}
\psi_a \otimes \psi_b + (\text{spectral tail})$.
The $W$-algebra generators
$W_s = \psi_s - (1/\Psi)\,{:}J\,W_{s-1}{:} - \cdots$ are
obtained by the nonlinear Miura inversion. The coefficient
$(\Psi-1)/\Psi$ on $J \otimes W_{s-1}$ arises as
$1 - 1/\Psi$: the Drinfeld
coproduct contributes $1$ from the cross-term
$\psi_1 \otimes \psi_{s-1}$, and the Miura inversion
subtracts $1/\Psi$ from the ${:}JW_{s-1}{:}$ correction.
Since both the cross-term and the Miura correction are
$s$-independent in structure, the coefficient $(\Psi-1)/\Psi$
is universal.
\end{proof}

\begin{remark}[Verification and coincidental agreement]
\label{rem:e-miura-verification}
The coefficient $(\Psi-1)/\Psi$ equals $1/\Psi$ at $\Psi = 2$
and $1/2$ at $\Psi = 2$: these coincide. This is a
coincidental agreement that masks the distinction between
$(\Psi-1)/\Psi$ and $1/\Psi$ at a single point. Verification
requires $\Psi$-values at three or more points (e.g.\
$\Psi = 1, 2, 3$). At $\Psi = 1$: $(\Psi-1)/\Psi = 0$
(the coproduct is an algebra homomorphism, consistent with the
commutative limit); $1/\Psi = 1$ (wrong). Concrete
verification runs through spins $2$--$6$.
\end{remark}

\subsection{Verlinde recovery}

\begin{proposition}[Verlinde formula from ordered chiral homology]
\label{prop:e-verlinde}
At integer level $k \ge 0$, the dimension of the ordered
chiral homology of $V_k(\mathfrak{sl}_2)$ at genus $g$ recovers
the Verlinde formula:
\begin{equation}\label{eq:e-verlinde}
Z_g
\;=\;
\sum_{j=0}^{k} S_{0j}^{2-2g},
\qquad
S_{0j}
= \sqrt{\frac{2}{n}}\,\sin\!\biggl(\frac{\pi(j+1)}{n}\biggr),
\qquad n = k + 2.
\end{equation}
Closed-form polynomials
$P_g(n) = n^{g-1}(n^2 - 1) \cdot R_{g-2}(n^2)$
are proved through genus $6$:
\begin{align}
Z_0 &= 1,\qquad Z_1 = n - 1,\qquad
Z_2 = \frac{n(n^2-1)}{6},\label{eq:e-verlinde-low}\\
Z_3 &= \frac{n^2(n^2-1)(n^2+11)}{180},\qquad
Z_4 = \frac{n^3(n^2-1)(2n^4+23n^2+191)}{7560}.
\label{eq:e-verlinde-high}
\end{align}
The leading coefficients are
$\mathrm{leading}(Z_g) \sim
\zeta(2g-2)/(2^{g-2}\pi^{2g-2})$, consistent with the
asymptotic growth of modular forms.
\end{proposition}

\begin{proof}
The Verlinde formula computes the dimension of conformal blocks
at genus $g$ and level $k$. The ordered chiral homology at
level $k$ is the factorisation homology of $V_k(\mathfrak{sl}_2)$
over the ordered configuration space on a genus-$g$ surface.
Handle attachment (sewing along a pair of pants) and separating
factorisation are verified: the pair-of-pants fusion product
(Definition~\ref{def:e-pants}) reproduces the Verlinde
fusion coefficients. The closed-form polynomials are obtained
by evaluating $\sum_{j=0}^{k} S_{0j}^{2-2g}$ using the
explicit $S$-matrix entries and summing by residue calculus.
The rational generating function
$G_n(x) = \sum_j 1/(1 - a_j x)$ with
$a_j = n/(2\sin^2(\pi j/n))$ produces all $Z_g$ as
coefficients.
\end{proof}

\begin{remark}[Conformal anomaly obstruction]
\label{rem:e-conformal-anomaly}
The obstruction to a constant (non-spectral) coproduct for the
Virasoro algebra is the conformal anomaly: the quartic-pole
OPE $T(z)\,T(w) \sim (c/2)/(z-w)^4 + \cdots$ forces
a primitive ansatz
$\Delta(T) = T \otimes 1 + 1 \otimes T + X$ to satisfy
$(c/2 + c/2)/(z-w)^4$ on the left but $(c/2)/(z-w)^4$ on
the right. The mismatch $c/2$ is the conformal anomaly,
quantified by $\kappa(\mathrm{Vir}_c) = c/2$. At $c = 0$,
the obstruction vanishes and a constant coproduct exists.
At $c \neq 0$, the spectral parameter is forced.
\end{remark}

% ================================================================
% SECTION 10: Three computations
% ================================================================

\section{Three computations}
\label{sec:computations}

\subsection{The programme}

The $\En$ operadic circle assigns to each chiral algebra
$\cA$ a package of data at each operadic level:
$\Eone$ (bar complex, $R$-matrix, Yangian),
$\Etwo$ (derived centre, Gerstenhaber bracket),
$\Ethree^{\mathrm{top}}$ (topologised bulk, BV structure),
together with the shadow tower
$(\kappa, S_4, \Delta, \mathrm{class})$ and the genus
anomaly $\mathrm{obs}_g$. The three standard families
(Heisenberg, affine Kac--Moody, Virasoro) exhibit the
full range of behaviour.

\subsection{Heisenberg $H_k$: class G}

\begin{proposition}[Complete $\En$ data for $H_k$]
\label{prop:e-heis-complete}
The Heisenberg algebra $H_k$ at level $k$ has the following
$\En$ data:
\begin{center}
\renewcommand{\arraystretch}{1.15}
\begin{tabular}{lll}
\textbf{Datum} & \textbf{Value} & \textbf{Source} \\
\hline
\textup{OPE} & $J(z)J(w) \sim k/(z-w)^2$ & defining \\
% AP126 trace-form check: r(z)|_{k=0} = 0. Y.
\textup{$r$-matrix} & $r^{\mathrm{Heis}}(z) = k/z$ &
 Example~\ref{ex:e-heis-einf} \\
\textup{$\kappa$} & $\kappa(H_k) = k$ &
 $\mathrm{av}(k/z) = k$ \\
\textup{Class} & G &
 $S_4 = 0$, $\Delta = 8k \cdot 0 = 0$ \\
\textup{$R$-matrix} & $S(z) = \id$ & abelian \\
\textup{Yangian} & cocommutative & abelian \\
\textup{Koszul dual} & $H_k^! = \Sym^{\mathrm{ch}}(V^*)$,
 $\kappa^! = -k$ & $K = k + (-k) = 0$ \\
\textup{$\ChirHoch^*$} & $\bC \oplus \bC[-1] \oplus \bC[-2]$ &
 Example~\ref{ex:e-heis-chirhoch} \\
\textup{$\Ethree^{\mathrm{top}}$} & trivial
 ($\Pthree$ bracket zero) & abelian CS \\
\textup{SC-formality} & SC-formal &
 class G $\Leftrightarrow$ SC-formal \\
\textup{Genus anomaly} & $\mathrm{obs}_g = k \cdot \lambda_g$
 \textup{(uniform-weight)} & $d_{\mathrm{fib}}^2 = k\,\omega_g$
\end{tabular}
\end{center}
\end{proposition}

\begin{proof}
The OPE is the defining data. The $r$-matrix $r(z) = k/z$
is the $d\log$-absorption of the double pole: the OPE
$J(z)J(w) \sim k/(z-w)^2$ gives propagator residue
$\Res_{z=w} d\log(z-w) \cdot k/(z-w)^2 = k/(z-w)$, hence
$r(z) = k/z$. At $k = 0$, $r(z) = 0$: the $r$-matrix
vanishes in the abelian limit. The averaging map gives
$\kappa = \mathrm{av}(k/z) = k$. Since $S_4 = 0$ (the
algebra has no quartic pole), $\Delta = 8k \cdot 0 = 0$
and the shadow tower terminates at degree $2$: class G.

The $R$-matrix $S(z) = \id$ because the OPE is
$\Sigma_k$-symmetric (abelian). The Yangian
$Y(\widehat{\mathrm{U}(1)})$ is the universal enveloping
algebra of the abelian Lie algebra, which is cocommutative.

The chiral Hochschild cohomology has total dimension $3$
(Example~\ref{ex:e-heis-chirhoch}). The $\Pthree$ bracket
on $C^*(H_k)$ is zero because $[J, J] = 0$ (abelian).
The $\Ethree^{\mathrm{top}}$ structure is trivial (all
operations are determined by the commutative cup product).

SC-formality holds because the shadow tower terminates:
class G implies SC-formal
(the discriminant $\Delta = 0$ forces all higher SC operations
to vanish).
\end{proof}

\subsection{Affine Kac--Moody $V_k(\fg)$: class L}

\begin{proposition}[Complete $\En$ data for $V_k(\fg)$]
\label{prop:e-km-complete}
The affine Kac--Moody algebra $V_k(\fg)$ at non-critical
level $k \neq -h^\vee$, with $\fg$ simple of dimension
$d = \dim(\fg)$, has the following $\En$ data:
\begin{center}
\renewcommand{\arraystretch}{1.15}
\begin{tabular}{lll}
\textbf{Datum} & \textbf{Value} & \textbf{Source} \\
\hline
\textup{OPE} &
 $J^a(z)J^b(w) \sim kB^{ab}/(z-w)^2 + f^{ab}{}_c J^c/(z-w)$
 & defining \\
% AP126 trace-form: r(z)|_{k=0} = 0 for trace-form convention. Y.
\textup{$r$-matrix} &
 $r^{\mathrm{KM}}(z) = k\,\Omega/z$
 \textup{(trace-form)} &
 $d\log$ absorption \\
\textup{$\kappa$} &
 $\kappa = d(k + h^\vee)/(2h^\vee)$ &
 $\mathrm{av}(r) + d/2$ \\
\textup{Class} & L &
 $S_4 = 0$, shadow depth $r = 3$ \\
\textup{$R$-matrix} & $S(z) = \id$ &
 $\Einf$-chiral \\
\textup{Yangian} &
 $Y(\hat{\fg})$, cocommutative &
 $\Einf$-chiral \\
\textup{Koszul dual} &
 $V_k(\fg)^!$, $K = 0$ & $K = \kappa + \kappa' = 0$ \\
\textup{$\ChirHoch^*$} &
 $\bC \oplus \fg[-1] \oplus \bC[-2]$ &
 Example~\ref{ex:e-km-chirhoch} \\
\textup{$\Ethree^{\mathrm{top}}$} &
 nontrivial, $\simeq$ CFG CS &
 Theorem~\ref{thm:e-e3-identification} \\
\textup{SC-formality} & not SC-formal & class L, $\ell_3 \neq 0$ \\
\textup{Genus anomaly} &
 $\mathrm{obs}_g = \kappa \cdot \lambda_g$
 \textup{(uniform-weight)} & curvature
\end{tabular}
\end{center}
The Koszul complementarity $\kappa + \kappa' = 0$ holds
for all simple $\fg$ and all $k$.
\end{proposition}

\begin{proof}
The $r$-matrix uses the trace-form convention:
$r(z) = k\,\Omega/z$, where $\Omega$ is the inverse Killing
form Casimir. At $k = 0$, $r(z) = 0$: the $r$-matrix
vanishes (abelian limit). The averaging map gives
$\mathrm{av}(r(z)) = k\dim(\fg)/(2h^\vee)$ (the double-pole
channel contribution); the full $\kappa$ includes the
Sugawara shift $\dim(\fg)/2$ from the simple-pole
self-contraction:
$\kappa = \dim(\fg)(k + h^\vee)/(2h^\vee)$.
At $k = -h^\vee$ (critical): $\kappa = 0$, the curvature
vanishes, and the bar complex is uncurved.

Class L: the OPE has a double pole and a simple pole, giving
shadow depth $r = 3$. Since $S_4 = 0$ (the OPE has no quartic
pole), $\Delta = 0$ and the shadow tower terminates at degree
$3$.

The $\Ethree^{\mathrm{top}}$ structure is nontrivial: the
$\Pthree$ bracket
$\{\phi_a, \phi_b\} = (k + h^\vee)\,f_{ab}{}^c\,\phi_c$
is the linearisation of the Chern--Simons interaction, and
the identification with the CFG algebra is
Theorem~\ref{thm:e-e3-identification}.

$\ChirHoch^1(V_k(\fg)) \cong \fg$ as a vector space, with
total dimension $\dim(\fg) + 2$.
\end{proof}

\begin{remark}[The Sugawara shift]
\label{rem:e-sugawara-shift}
For the Heisenberg (abelian), $\mathrm{av}(r(z)) = \kappa$
holds directly. For non-abelian Kac--Moody, the full
$\kappa$ includes the Sugawara shift:
$\kappa = \mathrm{av}(r(z)) + \dim(\fg)/2$. The term
$\dim(\fg)/2$ arises from the simple-pole self-contraction
through the adjoint Casimir eigenvalue $2h^\vee$. This is a
fundamental difference between abelian and non-abelian
averaging, and is the reason $\kappa(V_k(\fg)) \neq 0$ at
$k = 0$ (the non-abelian Lie bracket persists).
\end{remark}

\subsection{Virasoro $\mathrm{Vir}_c$: class M}

\begin{proposition}[Complete $\En$ data for $\mathrm{Vir}_c$]
\label{prop:e-vir-complete}
The Virasoro algebra $\mathrm{Vir}_c$ at central charge $c$
has the following $\En$ data:
\begin{center}
\renewcommand{\arraystretch}{1.15}
\begin{tabular}{lll}
\textbf{Datum} & \textbf{Value} & \textbf{Source} \\
\hline
\textup{OPE} &
 $T(z)T(w) \sim (c/2)/(z-w)^4 + 2T(w)/(z-w)^2 + T'(w)/(z-w)$
 & defining \\
\textup{$r$-matrix} &
 $r^{\mathrm{Vir}}(z) = (c/2)/z^3 + 2T/z$ &
 $d\log$ absorption \\
\textup{$\kappa$} &
 $\kappa(\mathrm{Vir}_c) = c/2$ &
 census \\
\textup{Class} & M &
 $S_4 \neq 0$, $d_{\mathrm{alg}} = \infty$ \\
\textup{$R$-matrix} & $S(z) = \id$ &
 $\Einf$-chiral \\
\textup{Yangian} &
 $Y^{\mathrm{dg}}(\mathrm{Vir}_c)$, dg-shifted &
 $d_{\mathrm{gen}} = 3$ \\
\textup{Koszul dual} &
 $\mathrm{Vir}_{26-c}$, $K = 13$ &
 self-dual at $c = 13$ \\
\textup{$\ChirHoch^*$} &
 concentrated in $\{0,1,2\}$, polynomial &
 Theorem~\ref{thm:e-theorem-H} \\
\textup{$\Ethree^{\mathrm{top}}$} &
 conjectural \textup{(}general topologisation\textup{)} &
 Conjecture~\ref{conj:e-topologisation-general} \\
\textup{SC-formality} & not SC-formal & class M \\
\textup{Genus anomaly} &
 $\mathrm{obs}_g = (c/2) \cdot \lambda_g$
 \textup{(uniform-weight)} & curvature
\end{tabular}
\end{center}
\end{proposition}

\begin{proof}
The $r$-matrix: the OPE has a quartic pole $(c/2)/(z-w)^4$,
a double pole $2T(w)/(z-w)^2$, and a simple pole
$T'(w)/(z-w)$. The $d\log$ propagator absorbs one pole order
(the $r$-matrix pole is the OPE pole minus $1$): the
quartic gives cubic $(c/2)/z^3$, the double gives simple
$2T/z$, and the simple pole contributes a regular term. Hence
$r(z) = (c/2)/z^3 + 2T/z$, not quartic.

The scalar $\kappa = c/2$ is the unique family for which
$\kappa$ equals half the central charge. The Koszul dual is
$\mathrm{Vir}_{26-c}$ with $\kappa' = (26-c)/2$, giving
$K = \kappa + \kappa' = c/2 + (26-c)/2 = 13$. Self-duality
at $c = 13$: $K = 13$ and $\kappa = \kappa' = 13/2$.

Class M: the OPE quartic pole gives $S_4 \neq 0$, hence
$\Delta = 8\kappa\,S_4 \neq 0$ and the shadow tower is
infinite. The generating depth is $d_{\mathrm{gen}} = 3$
($m_3$ generates all higher operations), but the algebraic
depth is $d_{\mathrm{alg}} = \infty$ (all $m_k$ are
nonzero). The algebra is not SC-formal.

The $\Ethree^{\mathrm{top}}$ structure for Virasoro is
conjectural: the holomorphic--topological theory is $3$d
gravity, and the BRST complex requires constructing the
gravitational path integral. This is the content of
Conjecture~\ref{conj:e-topologisation-general}.
\end{proof}

\begin{remark}[The $S_4$ invariant]
\label{rem:e-s4-virasoro}
The $S_4$ invariant of Virasoro is
$S_4 = 10/(c(5c + 22))$, with large-$c$ asymptotic
$S_4 \sim 2/c^2$. The discriminant
$\Delta = 8\kappa\,S_4 = 8 \cdot (c/2) \cdot 10/(c(5c+22))
= 40/(5c+22)$ is nonzero for all $c \neq \infty$.
\end{remark}

\begin{remark}[Generating depth vs algebraic depth]
\label{rem:e-depth-distinction}
The Virasoro illustrates the distinction between generating
depth and algebraic depth. The operation $m_3$ generates all
higher $m_k$ by recursive application
($d_{\mathrm{gen}} = 3$: finite), but every $m_k$ is nonzero
($d_{\mathrm{alg}} = \infty$: the tower never terminates).
The generating depth is a finiteness statement about the
recursion; the algebraic depth is a statement about
vanishing. These are different invariants: the Virasoro has
finite generating depth but infinite algebraic depth.
\end{remark}

\begin{remark}[Comparison table: G, L, M]
\label{rem:e-glm-comparison}
The three computations exhibit the escalation from G to L
to M:
\begin{center}
\renewcommand{\arraystretch}{1.15}
\begin{tabular}{llll}
& \textbf{Class G (Heis)} & \textbf{Class L (KM)} &
 \textbf{Class M (Vir)} \\
\hline
OPE poles & $2$ & $2$ (dp) + $1$ (sp) & $4+2+1$ \\
Shadow depth & $r = 2$ & $r = 3$ & $r = \infty$ \\
$\Delta$ & $0$ & $0$ & $\neq 0$ \\
$\Ethree$ & trivial & nontrivial (proved) &
 conjectural \\
SC-formal & yes & no & no \\
$\ChirHoch$ dim & $3$ & $\dim(\fg)+2$ & polynomial \\
$K$ & $0$ & $0$ & $13$
\end{tabular}
\end{center}
The passage from G to L introduces the non-abelian Lie
bracket and the non-trivial $\Pthree$ structure. The passage
from L to M introduces the quartic pole, the infinite
shadow tower, and the gravitational sector.
\end{remark}

% ================================================================
% SECTION 11: The closing conjecture
% ================================================================

\section{The closing conjecture}
\label{sec:closing-conjecture}

\subsection{Statement}

The $\En$ operadic circle~\eqref{eq:e-circle} traverses
$\Ethree \to \Etwo \to \Eone \to \Etwo \to \Ethree$.
Theorem~\ref{thm:e-e3-identification} establishes that the
output $\Ethree$ is isomorphic to the input for simple $\fg$.
The closing conjecture asks whether this holds in full
generality: not just for the deformation families, but for
all $\Eone$-chiral algebras with topologisable bulk.

\begin{conjecture}[Drinfeld centre equals derived chiral centre]
\label{conj:e-drinfeld-center-equals-bulk}
Let $\cA$ be an $\Eone$-chiral algebra
\textup{(}Notion~B, Definition~\ref{def:e-five-notions}\textup{)}
and $U_\cA = \cA \bowtie \cA^!$ the Drinfeld double.
Assume the topologisation hypotheses of
Section~\ref{sec:topologisation}, so that the bulk
$\Ethree^{\mathrm{top}}$-algebra and the derived chiral centre
are defined. Then there is an equivalence of
$\Ethree^{\mathrm{top}}$-algebras:
\begin{equation}\label{eq:e-closing}
Z(U_\cA)
\;\simeq\;
\cZ^{\mathrm{der}}_{\mathrm{ch}}(\cA),
\end{equation}
where $Z$ denotes the Drinfeld centre and
$\cZ^{\mathrm{der}}_{\mathrm{ch}}$ the derived chiral centre.
\end{conjecture}

\subsection{Three obstructions}

Three obstructions to a proof are identified.

\begin{enumerate}[label=\textup{(\Roman*)}]
\item \emph{Pointwise reduction for class M.}
 For class G and L algebras (Heisenberg, affine Kac--Moody),
 the Drinfeld centre $Z(U_\cA)$ can be computed stalk-wise:
 the centre at each point of the Ran space agrees with the
 chiral Hochschild cochains, and the global assembly is
 straightforward. For class M algebras (Virasoro, $\cW_N$),
 the $\Ainf$ corrections at degree $\ge 3$ are invisible to
 the stalk-wise centre. The global assembly requires a
 coherent family of higher homotopies that is not produced by
 the pointwise computation.

\item \emph{Strict factorisation for $\cA^!$.}
 The Drinfeld double $U_\cA = \cA \bowtie \cA^!$ requires
 the Koszul dual $\cA^!$ to carry a factorisation structure
 compatible with $\cA$. The Koszul dual is defined as an
 $(\SCchtop)^!$-algebra
 (Proposition~\ref{prop:e-koszul-dual-sc}): the closed colour
 is Lie, the open colour is Ass. The passage from the
 $(\SCchtop)^!$-algebra to a factorisation algebra on the Ran
 space requires careful completion in the pro-nilpotent
 topology. The completion is non-trivial for infinite-dimensional
 chiral algebras: the pro-nilpotent filtration on $\cA^!$ may
 not be compatible with the factorisation product.

\item \emph{Bar-cobar vs chiral Hochschild.}
 Bar-cobar inversion produces the algebra $\cA$ from the bar
 complex $\Barord(\cA)$:
 $\Omega(\Barord(\cA)) \simeq \cA$. This is inversion, not
 the chiral Hochschild construction. The compatibility between the
 bar-cobar involution and the chiral Hochschild functor (which
 computes the derived centre from the bar complex) is an
 independent statement. Concretely: the convolution
 $\Hom(\Barord(\cA), \cA)$ computes
 $C^\bullet_{\mathrm{ch}}(\cA, \cA)$, and this must be
 compared with the Drinfeld centre $Z(\cA\text{-mod})$.
 The comparison requires a Morita-type equivalence between
 the module category of $\cA$ and the comodule category of
 $\Barord(\cA)$, which is the bar-cobar equivalence of
 modules.
\end{enumerate}

\begin{theorem}[Resolution of any one obstruction suffices]
\label{thm:e-one-obstruction}
If any one of the three obstructions
\textup{(I)--(III)} is resolved, the
conjecture reduces to a comparison of two
$\Etwo$-categories, which is then a consequence of the
higher Deligne conjecture.
\end{theorem}

\begin{proof}
\emph{If (I) is resolved:} the Drinfeld centre is computed
stalk-wise for all classes, so $Z(U_\cA)$ is computed as a
chiral Hochschild cochain complex. The comparison with
$\cZ^{\mathrm{der}}_{\mathrm{ch}}(\cA)$ is the chiral
centre theorem (Theorem~\ref{thm:e-center-hochschild}).

\emph{If (II) is resolved:} the Drinfeld double $U_\cA$ is
a well-defined $\Eone$-chiral algebra, and
$Z(U_\cA)$ is a well-defined $\Etwo$-braided monoidal
category. The derived centre
$\cZ^{\mathrm{der}}_{\mathrm{ch}}(\cA)$ is also an
$\Etwo$-algebra by the chiral centre theorem. The comparison
of the two $\Etwo$-algebras reduces to matching their
$\Pthree$-structures via $\Ethree$ formality
(Theorem~\ref{thm:e-e3-formality}), and the one-dimensionality
of the deformation space (for simple $\fg$) forces agreement.

\emph{If (III) is resolved:} the bar-cobar equivalence of
modules produces a Morita equivalence between $\cA$-modules
and $\Barord(\cA)$-comodules. The Drinfeld centre of the
$\Eone$-monoidal category of $\cA$-modules is then computed
by the chiral Hochschild cochains of $\Barord(\cA)$, which is
$C^\bullet_{\mathrm{ch}}(\cA, \cA) =
\cZ^{\mathrm{der}}_{\mathrm{ch}}(\cA)$ by definition.

In each case, the resulting comparison is between two
$\Etwo$-algebras with matching $\Pthree$-structure. The
higher Deligne conjecture promotes both to $\Ethree$, and
the $\Ethree$ identification
(Theorem~\ref{thm:e-e3-identification}) establishes the
equivalence.
\end{proof}

\subsection{Evidence}

\begin{proposition}[Evidence for the closing conjecture]
\label{prop:e-evidence}
The following cases are verified:
\begin{enumerate}[label=\textup{(\roman*)}]
\item \emph{Heisenberg $H_k$, $k \neq 0$
 \textup{(}class G\textup{)}.}
 The Drinfeld double $U_{H_k} = H_k \bowtie H_k^!$ has
 centre $Z(U_{H_k}) = \bC^3$ (Proposition~\ref{prop:e-heis-complete}),
 matching the derived centre
 $\cZ^{\mathrm{der}}_{\mathrm{ch}}(H_k) = \bC^3$
 (Example~\ref{ex:e-heis-chirhoch}). All three obstructions
 are trivially resolved: pointwise reduction works (class G),
 $H_k^!$ is a free-field algebra (factorisation automatic),
 and bar-cobar inversion is a quasi-isomorphism (class G,
 Koszul).

\item \emph{Affine Kac--Moody $V_k(\fg)$, simple $\fg$,
 $k \neq -h^\vee$ \textup{(}class L\textup{)}.}
 The $\Ethree$ identification theorem
 (Theorem~\ref{thm:e-e3-identification}) establishes the
 closing at the level of formal deformation families. The
 closing conjecture asks for the identification as
 $\Ethree$-algebras, which requires matching the
 Drinfeld centre with the derived centre. The pointwise
 reduction obstruction (I) is resolved for class L (the
 shadow tower terminates at degree $3$, so the stalk-wise
 centre sees all the data). Obstruction (II) requires the
 Koszul dual $V_k(\fg)^!$ to carry factorisation structure;
 this is proved for $\fg = \mathfrak{sl}_2$.

\item \emph{Numerical verification for $H_k$.}
 The five invariants of the Drinfeld centre
 ($\Etwo$ dimension, $\Pthree$ bracket, cup product,
 $\Ethree$ operation, obstruction class) match at six
 levels $k = 1, 2, \ldots, 6$. The naive centre has
 dimension $1$; the derived centre has dimension $3$
 (the $\mathrm{Ext}^1$ and $\mathrm{Ext}^2$
 contributions are invisible to the commutant).
\end{enumerate}
\end{proposition}

\subsection{The holographic interpretation}

\begin{remark}[The circle as holographic duality]
\label{rem:e-holographic}
If the closing conjecture holds, the $\En$ operadic circle
is a mathematical formulation of holographic duality in $3$d
holomorphic--topological field theory: the bulk
$\Ethree^{\mathrm{top}}$-algebra determines the boundary
$\Etwo$-chiral algebra, its $\Eone$-bar complex, the
Drinfeld centre, and the derived centre, each recovering
the next. The passage from the ordered bar to the symmetric
bar is the operadic shadow of the passage from the boundary
to the bulk. The five theorems of modular Koszul duality
(A, B, C, D, H) are the $\Sigma_n$-invariant projections
of this holographic structure.
\end{remark}

\begin{remark}[The three obstructions as mathematical problems]
\label{rem:e-three-problems}
Each obstruction identifies a concrete mathematical problem:
\begin{enumerate}[label=\textup{(\Roman*)}]
\item The class~M pointwise reduction requires understanding
 $\Ainf$ corrections to the Drinfeld centre for algebras
 with infinite shadow tower. This is related to the
 chain-level topologisation problem.
\item The factorisation structure on $\cA^!$ requires
 constructing the Koszul dual as a genuine factorisation
 algebra, not merely an $(\SCchtop)^!$-algebra. This is
 related to the pro-nilpotent completion problem.
\item The bar-cobar vs chiral Hochschild comparison requires
 extending the Keller--Lef\`evre-Hasegawa theorem to the
 chiral setting: the derived Morita equivalence between
 $\cA$-modules and $\Barord(\cA)$-comodules must be proved
 in the chiral world.
\end{enumerate}
Resolution of any one of these would close the $\En$ circle
for all standard families.
\end{remark}

\begin{thebibliography}{99}
\bibitem{Arnold69} V.~I.~Arnold, \emph{The cohomology ring of the
 group of dyed braids}, Mat.\ Zametki \textbf{5} (1969), 227--231.

\bibitem{AF15} D.~Ayala and J.~Francis, \emph{Factorization homology of
 topological manifolds}, J.\ Topol.\ \textbf{8} (2015), 1045--1084.

\bibitem{BD04} A.~Beilinson and V.~Drinfeld, \emph{Chiral Algebras},
 AMS Colloquium Publications, vol.~51, 2004.

\bibitem{CFG25} K.~Costello, D.~Francis, and O.~Gwilliam,
 \emph{Perturbative quantization of Chern--Simons theory and the
 derived center}, in preparation, 2025.

\bibitem{Coh76} F.~R.~Cohen, \emph{The homology of
 $C_{n+1}$-spaces, $n \ge 0$}, in: The Homology of Iterated Loop
 Spaces, Lecture Notes in Math.\ \textbf{533}, Springer, 1976.

\bibitem{FM94} W.~Fulton and R.~MacPherson, \emph{A compactification
 of configuration spaces}, Ann.\ Math.\ \textbf{139} (1994), 183--225.

\bibitem{Francis2013} J.~Francis, \emph{The tangent complex and
 Hochschild cohomology of $\En$-rings}, Compos.\ Math.\ \textbf{149}
 (2013), 430--480.

\bibitem{Fresse17} B.~Fresse, \emph{Homotopy of Operads and
 Grothendieck--Teichm\"uller Groups}, Mathematical Surveys and
 Monographs, AMS, 2017.

\bibitem{FresseWillwacher20} B.~Fresse and T.~Willwacher,
 \emph{The intrinsic formality of $\En$-operads}, J.\ Eur.\ Math.\
 Soc.\ \textbf{22} (2020), 2047--2133.

\bibitem{GJ} E.~Getzler and J.~D.~S.~Jones, \emph{Operads, homotopy
 algebra and iterated integrals for double loop spaces}, preprint,
 1994.

\bibitem{GK98} E.~Getzler and M.~M.~Kapranov, \emph{Modular operads},
 Compos.\ Math.\ \textbf{110} (1998), 65--126.

\bibitem{HA} J.~Lurie, \emph{Higher Algebra}, 2017.

\bibitem{Idrissi22} N.~Idrissi, \emph{Real homotopy of configuration
 spaces}, Lecture Notes in Math.\ \textbf{2303}, Springer, 2022.

\bibitem{KhanZeng25} A.~Khan and Y.~Zeng, \emph{$\Ethree$-algebras
 from chiral algebras}, preprint, 2025.

\bibitem{Kon94} M.~Kontsevich, \emph{Feynman diagrams and
 low-dimensional topology}, in: First European Congress of
 Mathematics, Vol.~II, Birkh\"auser, 1994, 97--121.

\bibitem{Kon99} M.~Kontsevich, \emph{Operads and motives in deformation
 quantization}, Lett.\ Math.\ Phys.\ \textbf{48} (1999), 35--72.

\bibitem{Kon03} M.~Kontsevich, \emph{Deformation quantization of
 Poisson manifolds}, Lett.\ Math.\ Phys.\ \textbf{66} (2003), 157--216.

\bibitem{Livernet06} M.~Livernet, \emph{A note on the Lawrence--Sullivan
 interval and the Koszul property of the Swiss-cheese operad},
 preprint, 2006.

\bibitem{LV12} J.-L.~Loday and B.~Vallette, \emph{Algebraic Operads},
 Grundlehren der math.\ Wiss.\ \textbf{346}, Springer, 2012.

\bibitem{SalvatoreLongoni05} P.~Salvatore and R.~Longoni,
 \emph{Configuration spaces are not homotopy invariant},
 Topology \textbf{44} (2005), 375--380.

\bibitem{Tamarkin03} D.~Tamarkin, \emph{Formality of chain operad of
 little discs}, Lett.\ Math.\ Phys.\ \textbf{66} (2003), 65--72.

\bibitem{Totaro96} B.~Totaro, \emph{Configuration spaces of algebraic
 varieties}, Topology \textbf{35} (1996), 1057--1067.
\end{thebibliography}

\end{document}
